\chapter{O problema da desindustrialização na sociologia}

\section{Teoria do quê?}

Nas ciências sociais, raramente os problemas são prescritivos de uma
única abordagem teórica. Na verdade, distinto das ciências naturais, onde
a teoria assume quase sempre um papel de conhecimento acumulado sobre o tópico
em questão, cientistas sociais empregam a noção de teoria de forma mais ampla, abarcando tanto uma maneira particular de olhar e resolver seus problemas como enquanto uma orientação da própria identidade do pesquisador e do objetivo de sua pesquisa. É particularmente
inexistente para estes uma área de pesquisa onde "teoria" signifique apenas
um conjunto relativamente bem testado e reprodutível de hipóteses sobre dado
fenômeno \todouser{}{citar fonte}. 

O caráter \sout{multi-paradigmático}\aiflag{multiparadigmático} da sociologia \todouser{}{citar fonte},
como de diversas ciências ou áreas de estudo próximas, elimina
a possibilidade de \sout{se quer}\aiflag{sequer} um consenso sobre premissas epistemológicas.
Como consequência, seções de "teoria" acabam por quase sempre representar
um conjunto de premissas, autores e "linhas teóricas" aos quais os pesquisadores
tomam predileção, com frequência
mais em função de suas trajetórias institucionais do que propriamente pela
demanda de seus objetos de pesquisa.

\todoaiflag{}{SUGESTÃO DE TÓPICO (IA, 2026-10-19), pasta A.0: um exemplo concreto do caráter multiparadigmático é a oposição entre cultura/agência e estrutura, que atravessa a teoria social desde a polêmica Thompson × Althusser e foi nomeada por Hall como "dois paradigmas". As sínteses sociológicas (Bourdieu, Giddens, Archer) não encerraram a disputa, que segue viva (Fuchs 2023). Apoio: Hall 1980 e 1981; Thompson 1978; Anderson 1980; Giddens 1984; Archer 1995; Fuchs 2023.}

Essa característica implica em algumas consequências, boas e ruins. Do ponto
de vista positivo, a discussão teórica e metodológica nas ciências sociais
costuma adquirir um caráter mais robusto. Cientistas sociais pensam com mais
afinco sobre a natureza de seus objetos, a validade ou parcialidade de suas premissas,
o próprio lugar do pesquisador e as consequências de sua atividade \todouser{}{citar fontes}.

A contrapartida é que as escolhas teóricas com frequência enrijecem a capacidade
de observação, quando não \sout{pressupões}\aiflag{pressupõe} as próprias conclusões da investigação. Também
implica que aquilo que é chamado de teoria é com frequência muito distinto de um trabalho
para outro. Teorias de diversos níveis analíticos
se confundem e raramente o resultado de um trabalho é uma contribuição cumulativa
àquele campo de conhecimento -- ao menos não para além do grupo que
previamente já compartilhava da mesma "linha teórica".

Dada a multiplicidade de perspectivas válidas para abordar um problema
e a ausência de uma condição de ciência normal, talvez não seja possível solucionar
por inteiro esta problemática. O que se pode, entretanto, é sistematizar adequadamente
algumas respostas para facilitar ao leitor o entendimento do que, para o autor,
é teoria.

Assim, essa seção propõe esclarecer o que significa, no escopo deste trabalho
e no contexto da discussão sobre desindustrialização, o significado de teoria. Isso
significa esclarecer como instrumentalmente o papel desta é entendido na prática sociológica, quais
abordagens teóricas são consideradas mais frutíferas para o entendimento do objeto e
qual grau de conhecimento acumulado se tem sob o mesmo.

O segundo e o terceiro item são discutidos nas seções abaixo. Na seção 2.2 são
destrinchadas as principais linhas teóricas -- ou \emph{abordagens teóricas} -- \aiflag{que} são mobilizadas como explicativas, a saber:
a sociologia econômica, as vertentes institucionalistas e as discussões sobre desenvolvimento
e dependência. Já em 2.3 é resumido o que até o momento se produziu de mais relevante
sobre o fenômeno da desindustrialização no mundo e no Brasil.

Antes, porém, cabe explicar o uso de teoria neste trabalho. Ele é levemente diferente para os
dois itens descritos acima, justificando sua divisão em duas seções distintas. Dessa maneira, reflete certa separação entre uma teoria geral e uma
teoria substantiva \todouser{}{fonte}. A teoria geral, \sout{discutido}\aiflag{discutida} na seção 2.2, se refere \sout{as}\aiflag{às} grandes divisões internas a certa área disciplinar, que emergem, via
aglutinação ou diferenciação, pela confluência
temática mais direta -- sociologia do trabalho, sociologia do meio ambiente, sociologia da moral, etc --
ou pela semelhança mais indireta refletida em métodos, referências bibliográficas e conceitos compartilhados -- sociologia histórica,
sociologia econômica, sociologia política, etc.

Em ambos os casos essas grandes áreas \sout{sub disciplinares}\aiflag{subdisciplinares}
podem conter múltiplas subdivisões, as quais seria possível atribuir certa autonomia. Para aquelas reunidas por semelhança temática, a existência e crescimento de um número de trabalhos tematicamente mais próximo do objeto ou não parece justificar o enquadramento. No caso daquelas de união mais indireta, o que lhes diferencia e nos ajuda
a evitar um processo de \sout{hiper especificação}\aiflag{hiperespecificação} infinito
é seu grau de maturidade institucional. Esta é fruto de pesquisadores identificados com as mesmas que
articulam internamente maior coesão do que com aqueles "de fora". Esta coesão não pode então ser vista como natural, sendo geralmente construída intencionalmente
enquanto um projeto/agenda de pesquisa consciente por parte de pesquisadores -- representada pelos grandes volumes "definidores" desses campos, como "Método e visão
na Sociologia Histórica" \todouser{}{fonte} ou "Handbook of Economic Sociology" \todouser{}{fonte} para citar apenas dois exemplos relevantes a essa pesquisa.

Poderia ser argumentado que, em ambos os casos, esses grandes campos \sout{sub disciplinares}\aiflag{subdisciplinares} funcionam menos como uma teoria em estrito senso e mais
como "frames teóricos" \todouser{}{fonte} que orientam o conjunto de perguntas e os mecanismos interpretativos dos fenômenos.
Seu uso implica, portanto, tanto uma consequência da elaboração do problema, como também um posicionamento explícito do pesquisador
enquanto parte de \sout{algum(s)}\aiflag{algum(ns)} campo(s) de pesquisa, além de uma tentativa de delimitar áreas de atuação privilegiada. 
Considera-se que o termo "abordagem" é mais intuitivo na língua portuguesa
do que frame e por essa razão opta-se por definir estas como "abordagens teóricas". Legitimam e demandam, portanto,
"com o que" e "com quem" o pesquisador se propõe a trabalhar. O objeto deve ser adequável ao campo temático sob o qual aquela
abordagem reivindica legitimidade e os principais autores e conceitos precisam ser mobilizados, mesmo que criticamente ou em
discordância. Fazer uma "sociologia das emoções", por exemplo, implicaria passar obrigatoriamente por certos conceitos, autores, textos e métodos, mesmo que para rejeitá-los.

Não existindo nenhuma área que possa ser entendida como uma
"sociologia da desindustrialização", o mais próximo que se poderia chegar seria uma sociologia do desenvolvimento,
mas esta encontra-se em significativa paralisia nas últimas décadas \todouser{}{fonte}, o que tornaria seu uso
pouco produtivo. Opta-se assim pelo segundo tipo, menos temático e mais abstrato, o da abordagem teórica.
Entre as abordagens, diversas parecem ser importantes e contemplar o fenômeno da desindustrialização, 
mas há diferenças. Opta-se por um duplo critério, sabidamente subjetivo: deve possuir um conjunto de pesquisadores
atuantes e produtivos, para propiciar um diálogo vivo e uma contribuição maior da tese, e deve encapsular sem
significativos esforços a "desindustrialização" como um objeto ao qual tem legitimidade de discurso. Aquela que aparenta maior adequação 
e dispõe de um campo recente relativamente produtivo, é a sociologia econômica. Por obviedade isso implica
que se entende a sociologia econômica como uma abordagem teórica parte da disciplina sociologia e não como uma teoria específica,
uma disciplina a parte ou meramente um campo temático.

A seção seguinte desenvolverá esta justificativa e apontará, dentro desta escolha e em abordagens relacionadas, conceitos e autores que serão mobilizados, aqui por vezes tratados como vertentes internas a abordagem teórica mobilizada. Essas vertentes disputam espaço internamente e, no limite, poderiam adquirir autonomia suficiente para operar como abordagens teóricas próprias. Esse limite é evidentemente "embaçado", como qualquer categoria sociológica, e com frequência levantará disputas se \sout{esta}\aiflag{está} se fazendo uma "sociologia econômica de viés \sout{bourdiesiano}\aiflag{bourdieusiano}", uma "sociologia \sout{bourdiesiana}\aiflag{bourdieusiana} sobre a economia" ou uma "sociologia do mercado de arte \sout{bourdiesiana}\aiflag{bourdieusiana}", para limitar a poucas as possíveis classificações. Aqui não se argumenta ter feito uma taxonomia adequada da diversidade de qualquer disciplina ou sub-áreas, apenas que \emph{para fins desse trabalho} o que chamamos de teoria é a junção de uma abordagem teórica denominada sociologia econômica, com certas vertentes desta sendo privilegiadas, e uma teoria substantiva sobre desindustrialização, que abriga qualquer produção relevante sobre o tema que possa ser mobilizada enquanto conhecimento acumulado verificável.

Vale pontuar, outras duas abordagens teóricas identificáveis, a sociologia histórica e a sociologia do trabalho, foram observadas com maior frequência, 
quando na literatura sobre desindustrialização -- a teoria substantiva --, mas em última instância descartadas por menor dinamismo recente, 
no caso da primeira e por limitar demais as características relevantes do fenômeno a serem examinadas, no caso
da segunda. Nem por isso deixaram de ser mobilizadas em menor grau, juntamente de vertentes da sociologia política, a história econômica
e vertentes institucionalistas da ciência política e da economia, quando consideradas relevantes.

Já \aiflag{o} segundo item, discutido na seção 2.3, refere-se diretamente ao conjunto de trabalhos que investigam diretamente
o fenômeno da desindustrialização, sendo pertencentes a distintas abordagens teóricas. Aqui duas escolhas
foram feitas. Primeiramente, dividiu-se entre os trabalhos nacionais, sobre especificamente o
fenômeno brasileiro e aqueles internacionais. Segundo, dentro dos trabalhos nacionais a enorme
maioria encontrada se originou do campo da economia, o que dado o baixo número geral de trabalhos
sobre o tema, levou a sua incorporação. Para a literatura internacional, o número maior
de trabalhos permitiu -- e exigiu -- um recorte mais específico. Foram priorizadas então obras oriundas
da sociologia e da história, considerados os dois campos mais próximos desta tese. Esse segundo uso de "teoria" é, portanto, mais restrito e
mais próximo da descrição de \sout{teria}\aiflag{teoria} na metodologia científica tradicional. Ele permite elencar certo conjunto
de afirmações sobre o fenômeno ao qual seja possível atribuir certo grau de confiança e possam
em parte ser verificadas por meio da pesquisa empírica, assim fortalecendo ou enfraquecendo a confiança nelas.

\section{Os aportes da sociologia econômica}
Como discutido anteriormente, a abordagem teórica da 
sociologia econômica não se originou apenas do interesse
comum de parte dos sociólogos pelos temas econômicos.
Na verdade, a maioria dos autores considerados "clássicos" da disciplina lidava com a organização da economia como um de seus problemas centrais, seja Weber, Durkheim ou Marx \todouser{}{Swedberg}. Mas, assim como outros importantes contribuintes ao campo, como Veblen e Polanyi,
apenas seriam considerados fundadores ou figuras centrais deste em
retrospecto \todouser{}{fonte}. Se tomarmos os textos de Swedberg
\todouser{}{fonte}, \sout{Grannoveter}\aiflag{Granovetter} \todouser{}{fonte} ou Dobbin \todouser{}{fonte}, que operaram como organizadores conscientes da sociologia econômica nos anos \sout{1980 e 1980}\aiflag{1980 e 1990 (conferir)} em diante, foi o texto de Mark \sout{Grannoveter}\aiflag{Granovetter} "Estrutura Econômica e ação social" \todouser{}{fonte} que deu início a sociologia econômica como hoje estabelecida. Essa visão permanece mesmo em textos mais recentes, como Portes e Zelizer \todouser{}{fonte}

Foi nesse \sout{artivo}\aiflag{artigo} que \sout{Grannoveter}\aiflag{Granovetter} (re)introduziu o conceito organizador do que
seria a sociologia econômica, o \emph{enraizamento} \todouser{}{Fonte Swedberg}. O termo não era inédito, sendo mais identificado com a obra de Polanyi (fonte), onde trazia consigo sua contrapartida no
"desenraizamento". Ali, significava que a característica das relações econômicas na maior parte das sociedades era estarem condicionadas, limitadas
ou reguladas de alguma maneira pelas relações sociais entre os membros
\sout{daaquelas sociedade}\aiflag{daquelas sociedades} (fonte). O que Mauss (fonte) identifica na dádiva
em seus estudos sobre determinadas sociedades, um contraponto
\sout{a}\aiflag{à} maneira europeia moderna de pensar as \sout{relações econômica}\aiflag{relações econômicas}, 
Polanyi generaliza para um conjunto de alternativas de se lidar com
a organização econômica, variantes em acordo com as necessidades da
organização social (Polanyi). Necessidades essas que acabam deixadas
em segundo plano quando as relações econômicas adquirem autonomia
e primazia sobre as demais esferas da sociedade -- o desenraizamento --, com graves consequências para seus membros. Essas consequências engendram elas mesmas uma reação por parte da sociedade, que busca limitar o alcance e autonomia dos mecanismos de mercado e garantir a primazia de certas regras sociais (fonte idem).
As discussões sobre o conceito de enraizamento, particularmente o mecanismo de resposta da sociedade ao domínio do mercado \todouser{}{ver Burawoy} e as transformações que o termo sofreu na literatura e suas mutações ao longo do tempo \todouser{}{Cristiano e Rafael} ultrapassariam o escopo desta seção. \sout{Aqui o é importante}\aiflag{Aqui o importante é} esclarecer uma diferença fundamental, aquilo que a sociologia econômica entende como enraizamento, foi apenas em \emph{retrospecto} adequado ao conceito de Polanyi.

Isso pois, quando \sout{Grannoveter}\aiflag{Granovetter} (fonte) escreve, apesar da breve referência a Polanyi, seu uso do conceito é bastante distinto. Enquanto em Polanyi
a ideia de enraizamento opera num nível de abstração mais elevado, da
mesma forma que seu conceito de instituições, como será visto mais a frente, em \sout{Grannoveter}\aiflag{Granovetter} enraizamento é quase uma propriedade estrutural
de relações sociais empiricamente mapeáveis, ou "concretas" como ele se refere a elas \todouser{}{fonte}. Em \sout{Grannoveter}\aiflag{Granovetter}, dizer que algo
está \emph{enraizado} implica entender a posição relativa que o ator se
encontra frente a outros atores individuais ou coletivos e os fluxos
que se dão através de suas interações -- de recursos, informações ou
qualquer outro tipo. Tal perspectiva é oriunda da trajetória
intelectual do autor, que enfatizava uma análise estrutural e de redes, talvez
mais claramente explicitada em "A força dos Laços Fracos" \todouser{}{fonte} e em seu livro "Getting a Job" \todouser{}{fonte}.

Outra diferença e talvez a maior vantagem de \sout{Grannoveter}\aiflag{Granovetter} no uso do termo, é que para este as relações estudadas podem ser tanto
trocas econômicas como não-econômicas, sem comprometer sua capacidade de delimitar dada relação como \emph{enraizada}. A força de uma relação, seja dada pela frequência ou intensidade, é mais relevante que a mera característica privilegiada antecipadamente atribuída a um tipo de relação, por exemplo, a econômica. A dualidade, portanto reside muito mais na
diferenciação entre a escolha racional interna do indivíduo versus
a decisão condicionada externamente do que no conflito entre esferas sociais da vida e uma esfera econômica, por consequência "não social" como em Polanyi. \sout{Grannovetter}\aiflag{Granovetter} vai explicitar tal posição
na defesa de uma perspectiva sociológica que não seja nem sub e nem
"super socializada" (fonte). Se trata, portanto, de uma visão que prega um equilíbrio analítico entre a escolha racional e o contexto social -- a rede --,
sem necessariamente enfatizar a dinâmica entre esses dois elementos.

Já em Polanyi, as relações econômicas parecem ter um caráter qualitativamente distinto das demais. Ali, a dualidade -- ou dialética, apesar de Polanyi não se identificar explicitamente
com a tradição marxista -- está entre relações econômicas versus não econômicas, denominadas sociais. Essa qualificação particular das
relações econômicas tem um valor dúbio. Por um lado, de fato a especificidade das sociedades de mercado advindas do pós-século XVIII
parece justificar um tratamento distinto das relações econômicas, especialmente frente ao poder e influência que ganham. Ao mesmo tempo, ontologicamente não se sustenta essa diferenciação e acaba por encobrir o papel das relações econômicas
enquanto relações sociais, uma escolha talvez justamente para fugir do materialismo marxista ao qual Polanyi não queria ser atribuído.

Mas se em certo aspecto \sout{Grannovetter}\aiflag{Granovetter} avança, há um grande valor na
abordagem polanyiana, o que explica sua persistência e constantes 
"resgates" por autores da sociologia econômica. É que ao trabalhar
com a noção de enraizamento e desenraizamento, Polanyi estabelece uma
dualidade dinâmica e não estática. Não é meramente a necessidade de se optar por
uma perspectiva "equilibrada", como em \sout{Grannoveter}\aiflag{Granovetter}, mas uma que acrescente uma dimensão de tempo inata ao processo de transformação social. Que essa perspectiva tenha abrigado melhor trabalhos de cunho
histórico ou análises de nível macro-sociológico, não surpreende. Novamente, 
dinâmica polanyiana assume aqui um aspecto dialético bastante marxista, o que a torna útil a um grande conjunto de trabalhos
que seriam excessivamente constrangidos pelos requisitos do enraizamento \sout{grannoveteriano}\aiflag{granovetteriano}.

Dessa forma, \aiflag{a} influência de \sout{Grannoveter}\aiflag{Granovetter} na sociologia econômica levou a constituição de um grande conjunto de trabalhos identificados com o conceito de redes e preocupados com uma análise que \sout{contrabalancease}\aiflag{contrabalanceasse} as premissas utilitaristas dos economistas, sem necessariamente \sout{abraças}\aiflag{abraçar} uma super determinação das relações sociais. Já Polanyi acabou associado \sout{aqueles}\aiflag{àqueles} pesquisadores que tinham abordagens de fundo histórico \todouser{}{fonte} e também na maior parte do que seria uma vertente "institucionalista" \todouser{}{fonte}. Junto a Bourdieu e sua noção de campos, de alcance muito grande para ser inteiramente absorvido pela Sociologia econômica e que abrigou grande parte dos trabalhos de fundo culturalista da mesma, vão estabelecer um primeiro nível de divisão interna da sociologia econômica. Assim, pode-se definir ao menos três grandes vertentes analíticas da sociologia econômica, que representam formas similares, porém suficientemente distintas de interpretação da ideia de "enraizamento": redes, campos e instituições.

\subsection{Redes, Campos e Instituições}
Como dito acima, as três categorias que parecem organizar a forma
de pensamento da sociologia econômica são as redes, os campos e as
instituições. Pode-se dizer que, nessa ordem, representam graus crescentes de abstração das relações sociais e apoiam também crescentemente a dimensão temporal.

Assim, as redes são quem mais claramente delimita as relações sociais
a partir de sua representação empírica mensurável. Encontros, trocas,
co-participações... há algo que precisa \sout{se}\aiflag{ser} compartilhado entre dois indivíduos -- ou grupos -- e que estabelece entre eles uma relação. Essa relação pode frequentemente ser mensurada em intensidade ou frequência. O exemplo de seu desenho metodológico mais claro é o já citado "A Força dos Laços Fracos" (fonte) e sua aplicação "Getting a Job" (fonte). Ali, \sout{Grannoveter}\aiflag{Granovetter} mensura relações como trocas de informações visando a obtenção de uma vaga de emprego.

Mas há algo mais, pois o elemento central de \sout{Grannovetter}\aiflag{Granovetter}, assim como
de grande parte da análise de redes sociais (fonte) que ultrapassa
os cercados da sociologia econômica, é que a posição relativa entre os
elementos da rede pode importar até mais do que o número de relações ou sua intensidade.
Por exemplo, um elemento que intermedie a relação entre dois  outros elementos com grande número de relações cada um, um \emph{broker} na terminologia, pode ter
influência maior do que os \sout{mesmo}\aiflag{mesmos} na rede.

Da mesma forma, relações mais fracas, porém compartilhadas com membros
de \emph{clusters} distantes, agrupamentos de elementos, podem ser mais
eficientes na busca de informação, ou qualquer outro recurso disponível na rede, do que laços fortes estruturados
em grupos sociais bastante fechados em si. Esse, inclusive, é o principal argumento de "Getting a job" (fonte).

Por outro lado, a dimensão temporal acaba bastante negligenciada. Isso, pois os requerimentos para uma análise de redes propriamente \emph{\sout{grannoveteriana}\aiflag{granovetteriana}} já são bastante elevados, a coleta de dados deve ser exaustiva visto que poucos dados faltantes podem ter grande impacto nos resultados.
Da mesma forma as relações são muitas vezes informais e \sout{difíceis registro}\aiflag{difíceis de registrar}, recorrendo-se a métodos baseados na memória ou julgamento
dos indivíduos, que frequentemente omitem ou simplificam \sout{dinâmicas as reais}\aiflag{as dinâmicas reais}. Se considerarmos ainda os desafios matemáticos ligados a modelagem da rede, pois \sout{Grannovetter}\aiflag{Granovetter} enfatiza seu tratamento não como metáfora e sim estrutura modelável (fonte), não é de surpreender que poucos
trabalhos de fato consigam registros ou métodos adequados para modelar mudanças substanciais em períodos longos. Quando muito temos casos de antes e depois ou análises restritas a meios onde a coleta de dados seja por natureza facilitada, como redes digitais.

Muitos pesquisadores inclusive acabam por, na prática, utilizar o referencial \emph{\sout{grannoveteriano}\aiflag{granovetteriano}} sem de fato operacionalizá-lo. Não à toa
grande número de trabalhos emprega a "ideia" de redes sem efetivamente
mapeá-las ou fazê-lo apenas parcialmente (fonte). \sout{Dai}\aiflag{Daí} surgem as redes
como metáforas, que aceitam maior grau de abstração em prol da viabilidade de sua análise. São exemplos da vertente mais abstrata alguns estudos sobre redes de imigrantes (fonte) e de redes globais de produção (fonte). Vemos abordagens mais estruturais em estudos sobre gestão empresarial (fonte) ou interações em meio digital (fonte).

Um exemplo similar de maior abstração das relações sociais é o conceito de campo. Tornado altamente popular pela obra de Bourdieu (fonte), o campo representa justamente a delimitação de espaço -- geográfico, temporal, organizacional, etc -- onde se assume -- idealmente com base em alguma observação prévia, mas dificilmente mensurável -- um grau
de interação social, seja de cooperação ou disputa. Talvez justamente por essa flexibilidade e facilidade de operacionalização, campos são extremamente
frequentes como categoria analítica na sociologia e em outras ciências sociais. Como dito, a
vertente predominante é a \sout{bourdiesiana}\aiflag{bourdieusiana}, que coincide com grande ênfase
de seus praticantes em estudos sobre cultura ou grupos sociais, em
particular elites, tendo como exemplo os próprios trabalhos de Bourdieu (fontes). A ênfase aqui é em como aspectos da teoria \sout{bourdiesiana}\aiflag{bourdieusiana}, como a prática e o habitus, condicionam a ação de indivíduos em campos que são muitas vezes econômicos ou onde há forte prevalência de recursos econômicos em disputas: empresas, mercados, diretorias, elites econômicas, etc.

Assim como análises de rede, os estudos amparados no conceito de campo
tendem a ter dificuldades com a dinâmica temporal. Apesar do grande
esforço do próprio Bourdieu em capturar a dinâmica e transformação
de grandes campos em si (fonte), a maioria das análises tende
a ter dificuldade em rastrear um mesmo campo por longos períodos. \sout{Uma}\aiflag{Um} segundo obstáculo é a respeito da interação entre múltiplos campos e da potencial participação de indivíduos e grupos em diversos campos, já que o referencial \sout{bourdiesiano}\aiflag{bourdieusiano} centra-se em um campo por vez. Fligstein e \sout{Adams}\aiflag{McAdam} (fonte) fazem significativo esforço teórico para amparar análises de campos que considerem tanto transformações no tempo como a interação entre distintos campos. Fazem isso se 
afastando parcialmente dos conceitos \sout{bourdiesianos}\aiflag{bourdieusianos} e 
utilizando de repertórios oriundos da sociologia dos movimentos sociais (fonte). Um exemplo está na obra do próprio Fligstein (fonte) a respeito das transformações na governança corporativa das grandes empresas americanas e suas consequências.

Por fim, instituições avançam para um nível ainda maior de abstração -- das relações sociais --,
trabalhando não mais com as relações diretas ou os espaços, mas com
conjuntos de crenças cujo 'aspecto de regra' (fonte) molda a forma
de agir de grupos e indivíduos a partir de possíveis recompensas ou sanções. São a vertente mais frequente, mesmo que nem sempre totalmente explicitada, em trabalhos que lidem com 
níveis analíticos muito elevados ou longos horizontes temporais (fonte). Seu uso também é comum em diversas outras ciências sociais,
por exemplo, na economia e já era presente em trabalhos sociológicos
muito antes da formalização de uma vertente "institucionalista". Os
citados, Veblen e Polanyi são dois grandes exemplos, neles as instituições assumem um nível mais \sout{macrosociológico}\aiflag{macrossociológico}, respectivamente sobre classes sociais e a própria organização da economia. Já em autores ligados a sociologia econômica das organizações, como DiMaggio e FULANO (fonte), assumem um nível analítico intermediário. Instituições são formadoras de organizações e estas descritas como "ritos e mitos" que organizam a atividade econômica (fonte). Já economistas e parte dos cientistas políticos
tende a tomar as instituições mais explicitamente como as regras formais ou mesmo organizações reguladoras (fonte).

Aqui, em conjunto com a nova sociologia econômica (fonte nee), acreditamos nas instituições como regras formais e informais, onde há equilíbrio \sout{enter}\aiflag{entre} aspectos de racionalidade e costume. Instituições também nem sempre sobrevivem por serem a solução mais eficiente, mas dificilmente podem ser preservadas se não minimamente atenderem os requisitos funcionais, já que elas mesmas costumam \sout{demandas}\aiflag{demandar} recursos
para sua manutenção.

Também há predileção pelo uso de instituições ao nível analítico maior, como Polanyi faz em A Grande Transformação (fonte). Isso, pois também se lida com objetos cuja dimensão
geográfica e temporal é muito expansiva. A desindustrialização parece um caso claro onde, se é possível podemos olhar os impactos de dada legislação no fechamento de certa fábrica em dado ano, num horizonte mais largo precisamos de um enquadramento mais macro e em maior grau de abstração para capturar dinâmicas como conflitos
de classes e transformações no papel do Estado ao longo das décadas. \sout{Dai}\aiflag{Daí} nossas instituições privilegiadas serem "política cambial", "política industrial" e mesmo "características de \sout{frações classes}\aiflag{frações de classe}" e não "gabinete de planejamento" ou "práticas de gestão de investimentos".

De todas as três vertentes discutidas, o conceito de instituição também parece
aquele onde foi feito maior progresso em termos de sua análise de transformação no tempo. Verdade seja dita, grande parte dessa contribuição foi externa a sociologia econômica (fonte, fonte), mas hoje temos certo repertório para pensar não apenas como instituições
se mantém e se reproduzem, mas como nascem, morrem e as diferentes maneiras pelas quais se transformam (fonte). Novamente, isso torna uma vertente analítica bastante útil ao estudo de um fenômeno de longa duração como a desindustrialização.

\todoaiflag{}{SUGESTÃO DE TÓPICO (IA, 2026-10-18), pasta A.1: o papel do Estado na vertente institucional. O Estado como produtor e garantidor de instituições econômicas, e a autonomia inserida como forma de enraizamento do próprio Estado. Liga a escolha de instituições como política cambial e industrial ao debate sobre capacidade estatal. Apoio: Rueschemeyer e Evans 1985; Evans, Rueschemeyer e Stephens 1985; Evans 1995; Block e Evans 2005.}

Apesar das diferenças, uma semelhança das três vertentes, redes, campos e instituições, é a possibilidade de reenquadramento dos mesmos em função \aiflag{do} deslocamento do nível analítico. Todas as três acabam por operar a lógica de "bonecas russas", sejam campos dentro de campos, redes menores formando redes maiores, ou conjuntos de regras e normas distintas que podem ser lidos como instituições coerentes em si mesmas. Isso não muda que redes tendem a ser espacial e temporalmente mais "localizadas" ao passo que instituições possuem fronteiras menos espessas, com os campos sendo um caso intermediário.

E vale mencionar, \sout{tão pouco}\aiflag{tampouco} as três vertentes são mutuamente exclusivas ou possuem exclusividade sobre determinados objetos. A discussão aqui se deu em função da particularidade do problema dessa tese, a desindustrialização brasileira, mas outros problemas ou mesmo
um olhar distinto sobre este poderiam resultar em escolhas analíticas diferentes. Aliás, a escolha de objeto de estudo também é uma importante característica da sociologia econômica, novamente, se privilegiam tipos de objetos de estudo onde e através dos quais a noção de enraizamento produz seus efeitos. Esses objetos podem frequentemente ser lidos eles mesmos como redes, campos ou instituições. Abaixo discutimos cinco "tipos" de objetos mais característicos dessa literatura: organizações, mercados, grupos sociais, a própria história enquanto processo social e as culturas ou parte dessas.


\subsection{Os objetos onde esse enraizamento se desenvolve}
Num primeiro olhar a sociologia econômica poderia ser descrita como
abarcando toda e qualquer relação econômica. Não deixa de ser verdade
que é possível, mas a maioria dos trabalhos do campo acabam aderindo
a certos tipos de objetos. Assim, em níveis analíticos do menor para o maior:

Primeiramente, organizações. Grande parte da literatura do campo versa sobre 
organizações econômicas. Estas podem ser selecionadas por seu valor
histórico \todouser{}{citar}, \sout{pela}\aiflag{pelo} seu tipo de governança \todouser{}{citar} ou
por sua relevância e impacto em dado território \todouser{}{citar}. Em
geral, estudos sobre organizações econômicas nessa abordagem dão grande valor a estrutura da própria organização (fonte), a propriedade e governança da mesma (fonte), aos impactos sociais de suas atividades (fonte), a suas práticas gerenciais (fonte) ou a sua representatividade de algum caso \sout{suis generis}\aiflag{sui generis} (fonte). Parte da literatura de desindustrialização fala de casos particulares de organizações, normalmente do encerramento de suas atividades (fonte). Pode-se dizer tanto que as organizações estão enraizadas em dado contexto, por exemplo, o modelo de co-gestão do trabalho alemão (fonte), como que ações de membros ou grupos estão enraizados na própria organização, por exemplo, as práticas ambientais da Vale do Rio Doce (fonte). Em geral, transformações relevantes são aquelas que mudam ou o contexto no qual a organização está enraizada -- ex. uma nova legislação -- ou a própria organização enquanto fonte de enraizamento -- ex. a nova governança da Vale após sua privatização (fonte).

Em segundo lugar, mercados. Em geral, estudos sobre organizações conectam as mesmas ao seu contexto de mercado. Em alguns casos, o próprio mercado é o objeto chave, seja em seu processo de criação (fonte parpet) ou transformação (fonte). Como um \sout{intermediario}\aiflag{intermediário} temos também estudos sobre grupos ou clusters de empresas, ou ainda associações, que por vezes são entendidos como representativos de dado mercado (fonte) ou atores privilegiados em um deles (fonte). Os estudos de desindustrialização costumam enfatizar setores relevantes para economias nacionais ou locais que passam a ter significativa queda de emprego, como carvão (fonte) e aço (fonte). A própria definição de mercado evidentemente varia, sendo às vezes muito ampla, como em estudos sobre mercado financeiro (fonte), e por vezes bastante restrita, como em estudos sobre mercados de nicho. Enraizamento aqui costuma implicar nas preferências ou hábitos de um grupo consumidor particular (fonte) ou em determinadas estruturas de mercado que condicionem a operação das organizações (fonte).

Já os grupos sociais aparecem de duas maneiras. Alguns trabalhos enfatizam o papel ativo de dados grupos na articulação ou disputa por
certos mercados (fonte). É o caso de estudos sobre elites econômicas ou burocráticas (fonte). Já em outros casos os grupos sociais privilegiados são impactados pela atividade econômica de uma organização ou pela transformação de um mercado (fonte). São comuns exemplos onde a ênfase é na reação desses grupos (fonte) ou, ao contrário, que discutam a incapacidade de reação de um dado grupo (fonte). Muitas vezes estudos que enfatizam a particularidade de relações econômicas --impostas ou não -- a um grupo social fazem o recorte \sout{desde}\aiflag{desse} grupo a partir do gênero, como nas discussões sobre o papel da mulher no mercado de trabalho, de raça (fonte) ou classe (fonte). No caso da desindustrialização, a tendência observada foi de poucos trabalhos atribuindo agência ou capacidade de reação as comunidades afetadas e pouca problematização da agência das empresas, suas diretorias ou governos. Mas retornarei ao tema na seção seguinte.

Vale pontuar, nos três casos acima é comum a delimitação do objeto incluir um recorte geográfico explícito. Isso significa que empresas e mercados em geral estão "enraizados" em dado território (fonte) e que grupos sociais em geral são representativos de um dado território (fonte). Assim, temos trabalhos sobre o impacto de uma fábrica no sul do estado do Rio de Janeiro (fonte), sobre as mudanças na governança das empresas americanas (fonte) ou sobre \sout{as}\aiflag{a} invisibilização de trabalhadoras da \sout{frança}\aiflag{França} (fonte), para exemplificar. Enraizamento aqui é o enraizamento territorial, instituições, redes ou os \sout{habitos}\aiflag{hábitos} de dado grupo geograficamente localizado.

Curiosamente, a dimensão histórica é menos explicitamente articulada. Isso parece decorrer de grande parte da literatura operar sempre "no tempo presente", mas sem necessariamente problematizar o contexto social, político ou tecnológico. Como consequência, é difícil perceber num texto escrito em 1990 (fonte) uma diferença de contexto para um escrito em 2020 (fonte). O que é mais frequente é a contextualização de eventos recentes, entendidos, em geral, como elementos causais, mesmo que a relação causal não seja explicitada e nem possa ser inferida. Alguns estudos da sociologia do trabalho, que são tangenciais a sociologia econômica, tratam com mais especificidade as
mudanças de contexto político e tecnológico.

Nesse trabalho, dou mais peso a interpretação histórica, mesmo que os outros objetos entrecruzem. Nesse caso particular, enraizamento passa a ter grande relação com a noção de "trajetória de dependência" (fonte), já discutida no capítulo metodológico. Dizer que uma dada transformação está enraizada no tempo ou no processo histórico, implica assumir que escolhas e eventos deste possuem ramificações subsequentes que limitam a sua variabilidade, ao menos até a ocorrência de uma nova conjuntura \sout{critica}\aiflag{crítica}.

Por fim, a ideia de "cultura" como um conjunto de hábitos, valores ou ritos que perpassam determinado povo é um comum objeto de estudo da sociologia econômica. Tem-se desde estudos sobre demandas emocionais serem culturalmente atribuídas a dado gênero no mercado de trabalho (fonte), sobre mudanças no critério de uso e classificação do dinheiro (fonte) ou mesmo da valoração de acidentes ambientais entre distintas culturas (fonte). O que fica claro é que a dimensão cultural corta todos os casos acima e tem lugar bastante privilegiado
na sociologia econômica brasileira.

Dito isso, considero problemática a atribuição de relação causal entre uma "cultura" e dado fenômeno social, especialmente quando a própria cultura não pode ser facilmente recortada ou apresenta muitas mudanças ao longo do tempo. Em outras palavras, para estudar a relação entre cultura e desindustrialização, qual cultura deveria ser escolhida? A brasileira ou a dos territórios \sout{industrias}\aiflag{industriais}? A operária ou cultura das elites decisoras? Aquela das décadas de 1930, 1970 ou 1990? 

\sout{A demais}\aiflag{Ademais}, com \sout{frequẽncia}\aiflag{frequência} a sociologia econômica de viés culturalista
acaba por desenfatizar mecanismos de mercado ou fatores puramente econômicos. No objetivo de mostrar que não somos apenas maximizadores racionais, legítimo e correto, acaba por vezes assumindo um peso excessivo ao componente ideológico. Opera quase como uma estratégia de demarcação de espaço, voltando a discussão sobre a autonomia institucional das abordagens teóricas, onde os assuntos puramente econômicos são deixados aos economistas e caberia aos sociólogos ocupar os espaços deixados. Os mercados de nicho, os hábitos e os discursos. O risco é, ao meu ver, achar que entregadores de aplicativo o fazem não pela ausência de opções de trabalho, mas porque acreditam ser empreendedores, que os escravocratas buscavam
a superioridade racial antes do lucro e que as elites nacionais não
propõem o desenvolvimento meramente por serem "atrasadas". O padrão é uma relativa "fuga" das discussões de caráter material, para uma discussão "não materialista" da economia.

\todoaiflag{}{SUGESTÃO DE TÓPICO (IA, 2026-10-19), pastas A.0 e A.2: essa crítica tem precedentes. O balanço da virada cultural já perguntava se a redução do social ao discurso apagou instituições e economia (Bonnell e Hunt 1999); Jameson 1991 trata a cultura como lógica de uma etapa do capitalismo; Bourdieu ancora o gosto na posição de classe (Distinction). Para o exemplo dos entregadores de aplicativo: Burrell e Fourcade 2021 leem a precariedade das plataformas como relação de classe (elite do código × cibertariado), sustentada por propriedade e regulação.}

Vale dizer, de forma alguma a observação acima desqualifica as enormes contribuições dos estudos de viés culturalista. No Brasil, em particular sob a influência de Bourdieu, foram muito significativos os avanços a respeito de temas como a financeirização (fonte), para citar apenas um exemplo.

Alguns temas são mais frequentes na sociologia econômica brasileira. Financeirização, empreendedorismo, mercado de trabalho, mercados específicos -- inclusos ilegais -- práticas corporativas -- em particular de investimento. A financeirização é tipicamente vista por seu aspecto e impacto ideológico (fonte grun), não por uma transformação estrutural do modelo de acumulação como em Marx. 

\todoaiflag{}{NOTA (IA, 2026-11-01), pasta A.5: lista confirmada em parte. Lima e Higgins (2018, p. 13) enumeram quase os mesmos temas (gestão e sustentabilidade empresarial, empreendedorismo e construção de mercados, indústria automobilística e desenvolvimento regional, Estado, financeirização das classes populares, redes globais de produção); cite-os aqui. Mercados ilegais não aparecem nos mapas da pasta (nem em Vereta-Nahoum, 2017, pp. 34-35): precisa de fonte própria ou sai da lista. Sobre Grün: seu objeto é a dominação (elites, campo do poder, juros altos), lida por um método cognitivo e retórico, ``guerra cultural'' e ``convergência cognitiva'' (Vereta-Nahoum, 2017, pp. 38-41). ``Ideológico'' é termo seu, não dele; mais seguro: ``lida como dominação simbólica e cognitiva, não como mudança do regime de acumulação''. Ressalva do próprio balanço: essa dominação é ``matizada pela forte presença estatal'' (BNDES) (p. 41).}

Ela se conecta nessa dimensão ideológica/culturalisa a noção de empreendedorismo, que é em grande parte lido como uma ideologia ou como uma prática de construção de mercados. Em sua interpretação como ideologia, serve como parcial explicação ou efeito de um processo de precarização do trabalho (fonte). Porém, a questão aqui não é a disputa de classes per si inerente a extração do valor do trabalho, mas prioritariamente a disputa política em seu sentido mais estrito, partidário, onde a incapacidade de organização de partidos ou organizações favoráveis a pautas trabalhistas é lida como consequência da ideologia do empreendedorismo hegemonica (fonte). Assim, a disputa, apesar de política e econômica, é lida numa interpretação cultural. É nessa chave que costuma ser lida a transformação do mercado de trabalho. 

\todoaiflag{}{NOTA (IA, 2026-11-01), pasta A.5: nenhum texto da A.5 sustenta a leitura do empreendedorismo como ideologia que explica a fraqueza partidária dos trabalhadores; os dois (fonte) pedem obras de fora da pasta. O caso mais próximo na pasta é Sgorla, resenhado por Lima e Higgins (2018, pp. 16-17): a ``empresarialização de si'' de cervejeiros artesanais, via Sennett e Lahire, que trata de trajetórias, não de ideologia nem de partido. A passagem de ``disputa de classes'' para ``disputa partidária'' é a sua tese sobre a literatura: sem ao menos dois exemplos citados, fica atribuição genérica.}

\todoaiflag{}{NOTA (IA, 2026-11-01), bolsonarismo e precarização: a tese de que a precarização é lida sobretudo pela política institucional e pela ideologia, e não como conflito de classe, não se sustenta em termos gerais. Na sociologia do trabalho domina a leitura de classe (Braga, 2012; 2017; Filgueiras e Antunes, 2020; Druck, 2011; Marcelino e Galvão, 2020; Abílio, 2020), e Boito Jr. (2020) e Singer (2012; 2018) leem o bolsonarismo e a crise do lulismo como conflito de classe. A leitura pela subjetividade (Pinheiro-Machado e Scalco, 2020; Cesarino, 2022; Nunes, 2022; Rocha, Solano e Medeiros, 2021) recorre mais a Foucault e a Dardot e Laval do que a Boltanski, e rejeita o trabalhador ``convencido'' (Pinheiro-Machado e Scalco, 2021; Abílio, 2020). Ideologia e classe não se opõem: Antunes e Boito tratam a ideologia dentro da análise de classe. O contraste mais preciso é entre a ideologia como efeito das relações de classe e a cultura como causa autônoma. Obras conferidas por resumos e resenhas; confira as páginas antes de citar. Lista no Zotero, pasta 5 Leituras.}

Em geral, a precarização e liberalização são quase que efeitos dominantes, que refletem em vez de derivar, das transformações institucionais. Alguns poucos trabalhos abordam essas transformações do ponto de vista da classe e da disputa política (fonte santos e ramalho), mas a interpretação prioritaria parece uma de que o trabalhador teria sido "convencido" -- ou quase que doutrinado -- a adotar tal modelo e, diante dessa dominação ideológica, a resistência política seria fútil (fonte).

\todoaiflag{}{NOTA (IA, 2026-11-01), pasta A.5: (1) ``(fonte santos e ramalho)'': na biblioteca há Santos e Ramalho (2015, ``Estratégias corporativas e de relações de trabalho'') e Ramalho e Dos Santos (2018, sindicatos em redes globais de produção); confira qual. Vereta-Nahoum trata Ramalho e Santana (2006, Sul Fluminense) como literatura lateral (pp. 57-58), o que apoia ``alguns poucos''. (2) A tese do trabalhador ``convencido'' tem contraexemplo na própria pasta: Jardim (2009) e Grün (2009) mostram sindicalistas que passam a gerir fundos de pensão por mudança de habitus e estratégia de inclusão, não por doutrinação (Vereta-Nahoum, 2017, p. 40). E a linha de mercado de trabalho de maior destaque no balanço, Nadya Guimarães sobre busca de emprego, usa redes, não ideologia (p. 51). Sem fonte direta, ``doutrinado'' e ``fútil'' soam como caricatura: atenue ou cite casos.}

\todoaiflag{}{NOTA (IA, 2026-11-01), versão defensável deste argumento: em vez de dizer que a precarização é lida sobretudo pela política, delimite por campo. A leitura de classe da precarização está na sociologia do trabalho (Antunes, Braga, Druck, Galvão), que Vereta-Nahoum (2017, pp. 24-25, 57-58) trata como campo vizinho, e não como parte da sociologia econômica. Esta, quando estuda empreendedorismo, trabalha com sentidos e justificação (matriz Boltanski-Bourdieu, pp. 36-37; a ``empresarialização de si'' de Sgorla, em Lima e Higgins, 2018, pp. 16-17). Já os estudos do bolsonarismo explicam a guinada política dos trabalhadores sobretudo pela subjetividade. Papel da tese: trazer a leitura de classe para dentro da sociologia econômica.}

O empreendedorismo também é visto de uma outra maneira, pela consturção de mercados. Auqi, entretanto, apesar da influencia de Fligstein e sua perspectiva política, a chave principal é interpretar os mercados como "construções sociais". O ponto é que "sociais" nessa literatura não significa econômico, na verdade dizer que o mercado é uma consturção social implica quase dizer que o mercado é uma construção não econômica e sim cultural/ideológica (fonte).

\todoaiflag{}{NOTA (IA, 2026-11-01), pasta A.5: sustentado em parte. Vereta-Nahoum define a linha dominante como ``político-cultural'', lutas de poder pela definição das formas e práticas legítimas (pp. 34, 73), com Fligstein entre as matrizes (p. 37); e Raud (2007) lê em Bourdieu o Estado construindo oferta e demanda. Logo, ``social'' ali inclui poder e Estado, não só cultura. Crítica mais precisa e mais difícil de rebater: o social é político e simbólico, mas raramente estrutural (classe, acumulação, produção). Apoio direto: o balanço admite que ``são raros estudos no país [que] incluem materialidades das transações econômicas'' (p. 75), embora pense em artefatos e dispositivos, não em classe.}

Dai o número de estudos sobre pequenos mercados (como morango parpet) e mercados de nicho ou mesmo mercados ilícitos (fonte). A regra de ouro é encontrar o que a economia não explica e assim explicar via cultura, símbolos, tradições, críticas, morais ou racionalidades contrastantes a interpretação econômica dominante (fonte).

\todoaiflag{}{NOTA (IA, 2026-11-01), pasta A.5: o ``morango'' é Garcia-Parpet (1986), estudo francês (Fontaines-en-Sologne), que Vereta-Nahoum situa numa sociologia ``de matriz bourdiesiana'' (p. 54): use-o como influência, não como exemplo do campo brasileiro. Exemplo brasileiro de mercado de nicho: Sgorla, cerveja artesanal (Lima e Higgins, 2018, pp. 16-17). Para a ``regra de ouro'', aliado metodológico: Lima e Higgins criticam a redução de escala do estudo de caso único, que ``limita o poder inferencial da pesquisa'' (2018, p. 19). Ressalva: segundo Vereta-Nahoum, no Brasil o campo se formou ``menos em reação a explicações dos economistas'' e mais para renovar campos existentes (pp. 24-25), ainda que repita a retórica dos limites da economia; a ``regra de ouro'' vale sobretudo para a gênese estadunidense (Fourcade, 2007).}

Por fim, essa busca de "explicações alternativa" também domina diversos estudos sobre corporações e suas práticas (fonte). É especilamente importante a discussão sobre decisões de investimetno em face de uma lieteratura preocupada com o desenvolvimento regional (fonte ramalho). Estudos de caso em geral tentam capturar a implementação de inovações gerenciais (ramalho fabrica sonhos) que em geral são vistas como precursoras de precarização mais ampla do mercado de trabalho. Aqui há maior força da perspectiva política e econômica, em especial quando os estudos utilizam de modelos explicatóriso centrados na disputa e extração de valor, particularmente em redes globais de produção ou estudos sobre multinacionais. Mas é uma exceção a regra geral desenhada acima.

\todoaiflag{}{NOTA (IA, 2026-11-01), pasta A.5: bem sustentado, e é onde sua tese se encaixa. Práticas corporativas por ``explicações alternativas'': governança corporativa (Grün, 2005) e responsabilidade social via Boltanski e Chiapello (Carneiro, 2008, em Vereta-Nahoum, 2017, pp. 42-43). A exceção político-econômica: Dulci (2018, teoria da regulação, Sul Fluminense e Camaçari), Santos e Milanez (2018, redes globais de produção e Lukes) e Ramalho e Santana (2006), que o balanço trata como literatura lateral (pp. 57-58). ``(fonte ramalho)'' sobre investimento e desenvolvimento regional: candidato na biblioteca, Ramalho (2015, ``Indústria e desenvolvimento: efeitos da reinvenção...''). ``(ramalho fabrica sonhos)'': não achei título assim na biblioteca; talvez Ramalho (2006, ``Novas fábricas, velhas práticas''); confira. Sugestão da leitura A.5: em vez de só ``exceção à regra'', apresente-se como continuador dessa vertente, com a profundidade histórica e a escala que ela costuma não ter (Lima e Higgins, 2018, p. 19).}

\todoaiflag{}{NOTA (IA, 2026-11-01), pasta A.5, a partir da leitura integral de Vereta-Nahoum 2017 (``A Sociologia Econômica no Brasil: balanço de um campo jovem'', em Miceli e Martins, orgs., Sociologia brasileira hoje, pp. 15-86). (a) O que domina: uma abordagem ``político-cultural'', que entende as práticas econômicas como lutas de poder pela definição das formas organizacionais e práticas legítimas, ``uma sociologia política'' (pp. 34 e 73, seguindo Raud 2009). Suas três matrizes são a sociologia da crítica de Boltanski e Chiapello 1999, dita ``central'', o poder simbólico de Bourdieu 1977 e, explícita em menos estudos mas ``influente'', a abordagem político-cultural dos mercados de Fligstein (pp. 36-37). Zelizer orienta sobretudo outras linhas: dinheiro, circuitos e controvérsias morais (pp. 52, 56, 60-62). Redes (Granovetter) e performatividade (Callon) são ``residuais por aqui'' (pp. 34, 51, 54). Steiner aparece só como historiador do campo (pp. 18, 19, 25), não como matriz. (b) Formulação possível para este parágrafo e para ``Por fim, a ideia de cultura…'': na sociologia econômica brasileira predomina uma abordagem político-cultural (Vereta-Nahoum, 2017, pp. 34, 73), cujas matrizes são Boltanski e Chiapello, Bourdieu e, em menor medida, Fligstein (pp. 36-37). (c) Fonte para ``financeirização (fonte)'': a linha de Grün (2004; 2005; 2007; 2013) e Jardim 2009, que segue ``um esquema explicativo de inspiração bourdiesiana'' e lê a dominação financeira como ``guerra cultural'' e ``convergência cognitiva'' (pp. 37-41). Achado útil para o Contraponto: essa dominação é ``matizada pela forte presença estatal'' (BNDES) (p. 41). (d) Apoio e ressalvas ao argumento do autor: o próprio balanço admite que ``são raros estudos no país [que] incluem materialidades das transações econômicas'' (p. 75; também p. 43), mas ali materialidade significa artefatos e dispositivos, não classe ou produção. Recusa explicações estruturais: ``uma governamentalidade neoliberal, por si, não explica'' a governança corporativa e os fundos (p. 73), posição a contestar. Matiza a tese da demarcação: no Brasil o campo se formou ``menos em reação a explicações dos economistas'' e mais para renovar campos já existentes (empresas, desenvolvimento, trabalho) (pp. 24-25). A demarcação vale sobretudo para a gênese estadunidense (Fourcade 2007, p. 1017). Para o Brasil, o argumento mais preciso seria que as novas ferramentas atualizaram uma tradição sobre ``os dilemas do capitalismo brasileiro'' (pp. 35-36) deslocando seu centro das estruturas para os sentidos. (e) Lacuna a reivindicar: no mapa de cinco linhas do balanço (pp. 34-35), produção, política industrial e desindustrialização não aparecem. Detalhes na nota ``What dominates Brazilian economic sociology'' (Zotero, A.5).}

É curioso pois, em Bourdieu não se nega a perspectiva econõmica ou materialista (fonte), o habitus sendo justamente revelador das diferenças estruturais onde o dominante ocupa "uma posição na estrutura tal que a estrutura age em seu favor" \sout{(, 2005, p. 76)}\aiflag{(Bourdieu, 2005, p. 76)}. Além disso, próprio Bourdieu é crítico dessa separação onde aos sociólogos cabem os aspectos marginais da discussão econômica \sout{(bourdieu, 2005, p.8)}\aiflag{(Bourdieu, 2005, p. 84)}. Mas sua incorporação foi feita sobretudo pelo lado cultural (gosto, emoção, práticas, crenças). Creio que em muito pela pŕopria noção de campo levar ao estudo de disputas internas em grupos onde há similaridade de condições econômicas (acadêmico, de executivos, da arte) ou onde o modo de diferenciação mais vaorizado e eficiente frequentemtne é pelo status ou capital cultural. Mas um olhar para os diversos campos da sociedade mostra como eles se diferenciam, sobretudo, por suas características econômicas e políticas.

Foco em um campo -- leva a análise cultural
\todoaiflag{}{SUGESTÃO DE TÓPICO (IA, 2026-10-19, a pedido do autor), pastas A.1 e A.2: o uso culturalista do habitus. Parte da sociologia econômica toma de Bourdieu sobretudo o habitus pelo lado cultural (gosto, prática, emoções, crenças) e deixa de lado a estrutura que o produz: a distribuição desigual dos capitais, o campo como espaço de dominação e o Estado. Isso funciona como demarcação diante da economia (a cultura como objeto "próprio" da sociologia). O próprio Bourdieu recusa essa divisão de trabalho: critica a "distinção fundadora" que deixa aos sociólogos as curiosidades da economia (Bourdieu 2005, p. 84) e define o dominante como quem ocupa "uma posição na estrutura tal que a estrutura age em seu favor" (p. 76); o habitus econômico nasce das condições materiais (Argélia: Bourdieu 1979). A análise de classe de Bourdieu (espaço social, volume e composição do capital: Bourdieu 1985 e 1987; Wacquant 2013; Weininger 2005) seria o elo perdido. Apoio: Swedberg 2004, p. 12 (Bourdieu "muito mais estrutural") e 2011 (recepção restrita a A distinção); Lizardo 2011 (leitura culturalista "enganosa"); Raud 2007 (leitura brasileira do Bourdieu político e estatal). Exemplos brasileiros a avaliar: Grün 2004 e 2007; Vasques e Jardim 2021. Ressalvas: Burawoy 2010, p. 34 (o próprio Bourdieu examina as relações entre campos sobretudo pelo habitus) e Desan 2013 (Bourdieu toma o campo econômico como dado), que levam de volta ao Contraponto.}

Mas acredito que não devamos deixar a sociologia econômica ser uma sociologia daquilo que a economia não explica e nem que devamos fugir de uma disputa direta de espaço pela explicação dos fenômenos econômicos de ordem macrossociológica, a "economia" como sistema de relações sociais e não apenas dadas práticas ou mercados, o capitalismo como um arranjo institucional particular (FULANO), complexo e multi causal, mas explicável. Assim como a sociologia histórica teria predileção por "grandes questões" (\sout{sckopol}\aiflag{Skocpol}, fonte), creio ser necessária uma sociologia econômica que responda a grandes problemas -- o desenvolvimento, a luta de classes, o horizonte tecnológico da sociedade. E talvez a razão dessa carência  seja o relativo desaparecimento de uma importante vertente da sociologia.

Hall(fonte) critica propostas de retorno a perspectiva da economia política pelo risco de retomar antigos reducionismos economicistas. Mas no próprio autor e sua predileção pela perspectiva Gramsciana (fonte) como solução ao debate culturalista x estrutualista, acaba emergindo a importância da hegemonia. Esta por sua vez requer adequada consideção dos recursos materiais (fonte). Assim, é preciso diferenciar uma perspectiva materialista de uma perspectiva economicista. O materialismo aqui proposto é de ordem econômico-política, mas não ignora a dimensão histórica e da experiência que culturalistas como thompson(fonte) e Williams (fonte).

A cultura não está ausente. Ela é parte constitutiva do que é econômico e político. O conjunto de instituoções políticas e econômicas centrais a questão do desenvolvimento é constituido de e ele mesmo constitui uma dimensão cultural mais abstrata, que opera como discursos e visões sobre o desenvolvimento em si ou sobre o capitalismo. Dai uma análise institucionalista do capitalismo acabar por abordar um conjunto de características "culturais", como o ímpeto ao acumulo (streek). Esse diálogo entre níveis de abstração é justamente o que o método marxista propõe. Assim, o que se elimina é a distinção de uma dimensão "cultural" que opere apenas como discurso, símbolo ou ideologia, apartada de instituições político-econômicas que a reproduzam.

DESTACAR ASPECTO ENRAIZADO --> duplo enraizamento das instituições


\subsection{Contraponto: uma sociologia econômico-política}
Conclui na discussão acima com o argumento de que o referencial teórico e metologóco marxista precisa ser retomado como central a sociologia econômica. Mas é importante situar porque essa análise foi deixada em aprte de lado. QUe Marx é figura centrla da socioogia e ferquentemente citaod na sociologia econômica é fato defendido pela maioria dos programas da área (swedberg, portes, duina etc). Mas 

Em primeiro lugar
\todoaiflag{}{Tópico 1. Por que a sociologia econômica abandonou o nível macro: divisão de trabalho entre sociologia e economia herdada de Parsons, declínio da economia política marxista, fim da Guerra Fria e ida para as escolas de negócios. Apoio: entrevista com Schrank em González e Madariaga 2019; Burawoy 2003; Streeck 2012. Responde ao tópico sobre Marx como filho esquecido.}

Em segundo lugar, desde os anos 1980, em conjutno a "virada culturalista" (fonte), o marxismo perdeu grande parte de seu status no mundo ocidental. Isso se deveu em parte a novos referenciais que visavam lidar com problemas teóricos aos quais o marxismo não respondia adequadamente, mas também ao contexto político econômico. \todoaiflag{}{Tópico 10. Da polêmica marxista à sociologia, e o retorno da estrutura: Thompson × Althusser (anos 1970), sínteses sociológicas (Bourdieu, Giddens, Archer, anos 1980--90), virada cultural e seu balanço (Touraine 1984; Bonnell e Hunt 1999) e contracorrente macroestrutural (Panitch e Gindin 2012). Contextualiza o tópico sobre Marx como filho esquecido e o contexto dos anos 1980 em diante. Evidência de dentro do campo: na refundação dos anos 1980, "apelar ao trabalho de Marx tampouco parecia uma opção, uma vez que os tempos da sociologia radical haviam terminado" (Swedberg 2004, p. 16); no mesmo balanço, a própria definição do campo se apoia "no enfoque de Max Weber" (p. 7), e Marx aparece só como raiz distante, ao lado de Weber, Schumpeter e Polanyi (p. 7), e como referência perdida para o estudo da desigualdade (p. 22). Retorno recente da classe: Chibber 2022.}

Desde 1980, o crescimento econômico se estruturou em torno  da aceleração ed um processo de globalização combnado com adesão a um programa de liberalização da economia nos EUA e Europa Ocidental. A força desse modelo se refletiu na adesão de países Latino-Americanos a uma visãoa de desenvolvimento calcada na sua itnegração nesse novo "mercado globalizado", o que pode ser visto claramente nos processos de liberalização da economia dos anos 1980 e 1990. No caso latino-americano, esse processo ocorre juntamente a redemocratização de várias nações, o que ajuda a sedimentar uma perspectiva "otimista" com a dinâmica global.

Mesmo regiões onde o modelo econômico muito diferenciava do paradigma neoliberal, se é que ele existiu de fato em alguma economia central, o que era destacado eram sempre os processos de abertura comercial, investimento externo e entrada em novos mercados via exportações. Assim foi a trajetória de Japão, Coreias, dos tigres asiaticos e China (fonte). Não é que não houvesse distinção nos modelo econômicos, mas o que era politicamente valorizado desses, especialmente no discurso advindo das economias centrais, eram seus aspectos liberalizantes.

Associado ao fim da União Sovietica e no otimismo da possível integração dos estados do Leste europeu a união euripeia, que experienciava ainda significataivo sucesso, parecia que o modelo de mercado associado a democracia liberal ocidental, que constituem um só nexo, seria "hegemônico". Evidentemente, perspectivas que fossem centrlamente contrárias, como o marxismo e as diversas perspectivas centradas no imperialismo, sofreram forte retração. Não a toa muitos dos teóricos antes associados a sociologia do desenvolvimento e a teoria da dependência migraram para modelos menos centrados nas disputas de classe e sobre o Estado e mais orientados as disputas entre mercados, diferentes modos de organização dos mercados nacionais ou agentes econômicos. São exemplos os estudos sobre redes globais de valor, redes globais de produção, mercados financeiros, multinacionais ou mesmo sobre territórios. A quesao sai de confrotno ao modelo de acumulação, ou mesmo da possibilidade de desenvolvimento, e passa a ser sobre a forma de inserção no modelo.

A discussão sobre subdesenvolvimento enquatno condição local necessária ao modelo global passa a ser desvalorizada. Exemplos de sucesso "dentro do modelo" passam a ganhar atenção. A pergunta passa a ser como replicar o modelo de sucesso do Japão, da Coréia, de taiwan e mesmo da China. E em todos os casos a adesão ao modelo de mercado e a globalização. A desigualdade entre países passa a ser resultaod de boas ou más estratégias, não de dependência, e a desigualdade interna, a depender do olhar normativo escolhido, fruto ou de um regime meritocrático -- a desigualdade como boa e necessária -- ou de falta de justiça social -- a desigualdade como remediável pela justiça--, ambas interpretações morais e não estuturais. Se não chega a ser um retorno a teoria da modernização com suas fases pré-determinadas (fonte), especialmente por maior atenção as variações locais (fonte), reduz significativamente a discussão da desigualdade enquanto característica relacional entre países, típica das vertentes marxistas (fonte).

Mas há boas razões para crer no retorno da economia política, ou de uma sociologia econômico-política. Desde a crise de 2008, significativa parte do otimismo econômico com os processos de globalização se perdeu. Ao mesmo tempo, sua expressão política se encontra desafiada pelo retorno de movimentos anti-liberais. O pêndulo Polanyiano retorna contra o mercado, mas novamente o contexto de crise exige não apenas uma resposta da sociedade civil, como propusera Burawoy (fonte), mas uma disputa pelo Estado e entre estados. \todoaiflag{}{Tópico 8. Marxismo sociológico: a sociedade ativa de Polanyi e a sociedade civil de Gramsci como resposta à mercantilização. Faz a ponte de volta à seção 2.2 (enraizamento). Apoio: Burawoy 2003.}

Na dimensão geopolítica, o crescimento das economias do leste asiático passa não a ser uma coroação do processo de globalização, mas um desafio ao mesmo. Isso pois essa nova etapa do processo de acumulação capitalista centrada no leste asiático (arrighi) coloca em rota de colisão o novo e antigo centros, exigindo das potencias antigas o exercício cada vez mais aberto do poder militar e econômico de forma a manter suas áreas de influencia.

Assim, um retorno a Marx pode contribuir em diversas maneira. Primeiro,  recupera a violência não simbolica como fator explicativo das transformações (fonte). \todoaiflag{}{Tópico 12. A base material do capitalismo digital: mesmo leituras culturalistas da elite do código mostram propriedade dos meios técnicos e extração de trabalho (Burrell e Fourcade 2021). As disputas atuais por datacenters, terra, energia e semicondutores subordinam a dimensão intelectual (software, modelos) a um poder material tradicional e retomam trabalho e acumulação primitiva. Liga-se aos tópicos sobre poder não simbólico e disputa geopolítica. (A partir da reflexão do autor sobre Burrell e Fourcade.)}

Segundo, reintroduz o papel das classes sociais e dos conflitos entre e intra classes no estabelecimento de modelos de desenvolvimento estáveis. \todoaiflag{}{Tópico 11. O debate Miliband × Poulantzas como versão do dilema estrutura × agência na teoria do Estado: elites e vínculos (Miliband 1969, 1970 e 1973) contra posição estrutural (Poulantzas 1969 e 1973), com Block 1977 e Therborn 1978 como saídas. Cuidado útil para a "defesa de retomada estruturalista": a crítica de Miliband ao superdeterminismo estrutural.}

Também conecta a disputa geopolítica global as disputas locais que emergem com os muitos movimentos "anti-sistema" (ver euromaiden, 2013 no Brasil, hong kong). \todoaiflag{}{Tópico 4. Poder, coerção e Estado na formação do capitalismo e na trajetória dos países centrais. Apoio: Marx, acumulação primitiva; Block 2018; Chang 2002 sobre o protecionismo histórico dos países ricos.}

\todoaiflag{}{Tópico 5. Hegemonia, ciclos sistêmicos e a ascensão da China como pano de fundo da disputa geopolítica atual. Apoio: Arrighi, O longo século XX e Adam Smith em Pequim; Wallerstein; Frank, ReOrient. Ressalva: o sistema-mundo é visto como paradigma em declínio (Portes 2015; Viterna e Robertson 2015).}

Por fim, ajuda a retomar a centralidade do Estado nas discussões econômicas, coerente com a nova força de modelos centrados na ação estatal, sejam eles descritos como desenvolvimentismo, socialismo de mercado, capitalismo de Estado ou qualquer outra denominação. Estado e economia como esferas mutuamente constitutivas: crítica da crítica à sociologia estadocêntrica. Apoio: Skocpol em Evans, Rueschemeyer e Skocpol 1985; Block e Evans 2005; Evans 1995.

Mas retornar à Marx não significa não incorporar novas capacidade teóricas que emergiram. Como dito anteriormente, aqui a que mais interessa são os diversos desenvolvimentos entendidos sob a alcunha de "institucionalismo". Isso porque a noção de instituição é suficientemente flexível para permitir o trabalho e relação entre os níveis de abstração distintos. Também contribui pois as instituições, por seu caráter de crença e norma, tornam explícita a incorporação do elemento culturalista ou ideológico. Assim, falar de instituições econômcias e políticas é também falar dos discursos, símbolos e ideias, sem com isso perder a capacidade de diferenciação dos diversos elementos constitutivos do todo, que era contribuição do estruturalismo (hall).


\todoaiflag{}{SUGESTÕES DE TÓPICOS (IA, 2026-10-19), rodada Cultura e Estrutura (pastas A.0 e A.2), a partir das notas Ficha IA no Zotero. Mesma regra: manter, fundir ou descartar.}


\subsection{Institucionalismo(s), desenvolvimento e a dependência}
Falar de silver --> pendulo e etapas

INSTRITUCIOANLISMO


O que é o capitalismo institucional

\todoaiflag{}{Tópico 3. O capitalismo como ordem social e arranjo institucional historicamente específico, e não só como economia: ponte entre sociologia econômica e economia política. Apoio: Streeck 2011 e 2012; Beckert e Streeck 2008.}


\todoaiflag{}{Tópico 9. Classes, Estado e revolução burguesa no Brasil como caso do capitalismo dependente: a exceção brasileira mantida na A.2. Apoio: Fernandes 1975.}

Necessidade de um contraponto institucionalista (polanyi → evans)
Institucionalismo(s)
Sociologia econômica (histórica-comparativa e institucionalista) do
desenvolvimento

institucionalismo (economico, cultural e político)
Diferentes abordagens mono institucionalista ou apenas relacionando
cultura e economia (marx, weber, schumpeter, veblen, elias)

-- nota pessoal (2026-10-31, a partir de Hall 1980): talvez a abstração possa partir das próprias instituições econômicas e políticas, que carregam em si a cultura/ideologia e ao mesmo tempo estão dialeticamente imersas nela. Assim, noções ideológicas como desenvolvimento, nação, povo e classe seriam fundamentais para a construção das instituições políticas e econômicas (ao mesmo tempo que são definidas por elas), inclusive da ``classe para si''.

\todoaiflag{}{AVALIAÇÃO E REFERÊNCIAS (IA, 2026-10-31) para a nota acima. (a) Consistência: a síntese é coerente e evita os dois polos de Hall (1980): não trata a cultura como reflexo do econômico nem como totalidade que dissolve as instituições, e dá à ideologia um lugar concreto, as instituições. Tem antecedentes fortes, o que a torna defensável, mas pede que se diga o que ela acrescenta a eles. (b) Antecedentes diretos: Gramsci, ``bloco histórico'' (unidade entre estrutura e superestruturas) e ``Estado integral'' (sociedade política mais sociedade civil), Cadernos do cárcere; Althusser 1970, a ideologia tem existência material, existe sempre num aparelho e em suas práticas (conferir a citação); Sewell 1992 (AJS, no Zotero A.0), estruturas como esquemas e recursos que se sustentam mutuamente; Hall 1980, pp. 68-69, articulação; Polanyi 1944 (no Zotero), a sociedade de mercado como projeto institucional guiado pelo ``credo liberal'' (``o laissez-faire foi planejado''). (c) Para cada noção: desenvolvimento, Sikkink 1991 (Ideas and Institutions: Developmentalism in Brazil and Argentina), a ideia desenvolvimentista construindo instituições estatais no Brasil, e Bielschowsky 1988 (Pensamento econômico brasileiro: o ciclo ideológico do desenvolvimentismo); nação e povo, Benedict Anderson 1983 (Comunidades imaginadas); classe para si, Thompson 1963 (no Zotero; prefácio: a classe ``acontece'' na experiência e na luta) e Przeworski 1977 (Politics and Society), a formação de classe como efeito de lutas de partidos, sindicatos e Estado, e não passagem automática da posição à consciência; ponte institucionalista, Hall e Taylor 1996 (no Zotero; institucionalismo sociológico) e Blyth 2002 (ideias na mudança institucional). (d) Riscos: (1) se toda instituição carrega cultura e toda cultura está em instituições, o argumento volta à ``totalidade expressiva'' que Hall critica (1980, p. 68); é preciso dizer o que as instituições têm de próprio: recursos materiais, coerção, regras (Gramsci: hegemonia ``protegida pela couraça da coerção'', conferir). (2) A via de mão dupla (as noções constroem as instituições e são definidas por elas) é o ponto mais forte da síntese, mas também o que a expõe à circularidade; ``dialeticamente imersas'' pede mecanismo e tempo: quando a ideologia molda instituições e quando instituições já estabelecidas produzem ideologia; o dualismo analítico de Archer 1995 (no Zotero A.0) ajuda a separar as fases. (3) A ideologia pode construir a classe para si (Thompson, Przeworski), mas a classe em si não deve virar efeito do discurso; Chibber 2022 (no Zotero A.0) é o contraponto obrigatório. (e) Formulação sugerida: instituições como cristalizações materiais de projetos ideológicos vitoriosos, que passam a constranger as disputas ideológicas seguintes. No caso brasileiro: o desenvolvimentismo como ideologia que construiu o Estado industrializante (BNDE, estatais, planejamento) e cujo desmonte acompanha a desindustrialização. Fora do Zotero: Sikkink, Bielschowsky 1988, B. Anderson, Przeworski, Blyth, Althusser 1970.}

Polanyi como institucionalismo econômico e político


\todoaiflag{}{SUGESTÕES DE TÓPICOS (IA, 2026-10-18) para Institucionalismo(s), desenvolvimento e dependência (pasta A.3), a partir da seleção bibliográfica e das notas Ficha IA no Zotero. O esboço acima é do autor; os tópicos abaixo são sugestões a manter, fundir ou descartar.}

\todoaiflag{}{Tópico 1. Da dependência à virada institucional: por que as instituições passaram a ocupar o lugar explicativo da dependência, e o que se perdeu nessa troca. Apoio: Portes 2015; Viterna e Robertson 2015; Schrank 2023.}

\todoaiflag{}{Tópico 2. Três versões da dependência: desenvolvimento do subdesenvolvimento (Frank 1966 e 1967), dependência associada (Cardoso e Faletto 1970) e tríplice aliança entre Estado, multinacionais e capital local (Evans 1979). Balanços: Heller, Rueschemeyer e Snyder 2009; Bastos 2025; Giorgini e Iumatti 2025.}

\todoaiflag{}{Tópico 3. Estruturalismo e subdesenvolvimento como processo autônomo, e não como etapa: Furtado 1959, 1961 e 1974; Caio Prado 1968; Mello 1982 sobre o capitalismo tardio.}

\todoaiflag{}{Tópico 4. Classes sociais e dependência: a leitura de Florestan Fernandes (1968, 1973, 1975) como contraponto sociológico ao estruturalismo econômico.}

\todoaiflag{}{Tópico 5. Polanyi e Evans como eixo do institucionalismo político: Estado inserido, mercadorias fictícias e regulação. Apoio: Evans 1995; Block e Evans 2005; Polanyi 1944.}

\todoaiflag{}{Tópico 6. Mudança institucional ao longo do tempo, útil para um fenômeno de longa duração como a desindustrialização: Streeck e Thelen 2005; Mahoney e Thelen 2010; Duina 2011 para as definições de instituição.}

\todoaiflag{}{Tópico 7. Estado e mercado no sistema mundial: condições da intervenção estatal eficaz e o lugar dos Estados periféricos. Apoio: Rueschemeyer e Evans, The State and Economic Transformation (1985); Evans, Rueschemeyer e Stephens (orgs.), States versus Markets in the World System (1985); Heller, Rueschemeyer e Snyder 2009.}


\todoaiflag{}{BIBLIOGRAFIA (IA, 2026-10-29, pasta A.3) para "marxismo perde força". (a) Declínio da dependência a partir dos anos 1980 por razões políticas (fim do ciclo cubano; Katz e Packenham): Giorgini e Iumatti 2025, p. 3. (b) "Toda época produz seu próprio marxismo"; o marxismo depois de 1989: Burawoy 2003, p. 194. (c) Ocaso da sociologia do desenvolvimento brasileira nos anos 1980: Santos 2016, p. 296; Leme 2015, pp. 495-496. (d) Volta recente do tema no Norte: Bastos 2025, p. 2; Kvangraven 2021; Ellner 2024. A virada cultural já está na 2.1/2.2.2 (Chibber 2022): remeter, não repetir.}


O caso do sul global (institucionalismo e dependência)
Imperialismo
--tradição de literatura marxista
--argumentar que a tradição marxista de pensar imperialismo continua relevante na interpretação do capitalismo global
--citar contexto atual de disputa por mineráris e pressão geopolítica do governo trump sobre a américa Latina
--reforçar disputas políticas nacionais como disputas de modelos de desenvolvimento mais ou menos dependentes
--apesar de marx e polanyi permitirem um diálogo institutcionalista importante (silver), são expressões vindas do Norte


Teoria da dependência e desenvolvimento dependente
--desenvolvimento dependente como característica de nosso desenvolvimento
--dependencia como frame teórico e não teoria (texto reuschmayer)
--teoria da dependência aponta as relações chave entre a dinâmica interna e externa
--assim funciona como perspectiva sistema mundial (arrighi, wallerstein)
--ao mesmo tempo destaca disputa interna e luta de classes

\todoaiflag{}{BIBLIOGRAFIA (IA, 2026-10-29, pasta A.3) para "Teoria da dependência e desenvolvimento dependente". (a) Desenvolvimento dependente: Cardoso e Faletto 1970; Evans 1979, pp. 9-11 (semiperiferia; o desenvolvimento dependente não elimina a dependência); Mello 1982. (b) "Texto reuschmayer": a formulação literal "theoretical frame" está em Mahoney e Rodríguez-Franco 2018, p. 22 (já anotado por você); somar Heller, Rueschemeyer e Snyder 2009, Palma 1978 e Cardoso 1977 (a dependência como metodologia de situações concretas); Kvangraven 2021 fala em "programa de pesquisa". (c) Interno e externo: Dos Santos 1970, p. 231; Fernandes 1968 (dupla articulação). (d) Sistema-mundo: Wallerstein; Arrighi 1994; Mahoney e Rodríguez-Franco 2018, p. 23 (sistema-mundo como derivação da dependência); Frank como cofundador (Giorgini e Iumatti 2025, p. 5); Brenner 1977 como crítica marxista (fora do Zotero). (e) Classes: Fernandes 1968, 1973 e 1975; Corsi 2022; Marini 1973; Bambirra 1978. (f) Duas vertentes, econômica e política: Bastos 2025, pp. 1-3.}

DESENVOLVIMENTO surge para pensar o problema do Sul Global, não do Norte, mas nasce orientado pelo norte (o norte não se entende como desenvolvido sem o Sul existir como subdesendolvido --> decorrencia do imperialismo e da perspectiva relacional)
--como o sul global pensa/repensa suas relações de dependência e a economia política?
--pensar em trazer pensadores chineses

\todoaiflag{}{BIBLIOGRAFIA (IA, 2026-10-29, pasta A.3) para "Imperialismo". (a) Tradição marxista vista da periferia: Baran 1957; a dependência como "releitura do conceito e da história do imperialismo desde o Terceiro Mundo" (Bastos 2025, p. 3). (b) Relevância atual: Harvey 2003 (acumulação por espoliação); Katz 2021; fora do Zotero: Patnaik e Patnaik 2021, Kohli 2020. (c) Minerais e Trump: usar imprensa de 2026 só em rodapé (acordo de minerais críticos com preço mínimo; reunião Lula-Trump em maio com terras raras e processamento local; compra da Serra Verde); a ponte teórica é o neoextrativismo (Svampa 2019, fora do Zotero) e Bielschowsky 2012 como contraponto (recursos naturais como "frente de expansão", p. 730). (d) Disputas nacionais como modelos de desenvolvimento: dependência como projeto das classes dominantes (Fernandes 1973; Corsi 2022, p. 241); tríplice aliança (Evans 1979, p. 100). (e) Marx e Polanyi via Silver: Silver e Arrighi 2003 (Politics e Society) é o texto exato; Silver 2003 cap. 1 (agitações marxiana e polanyiana); Fraser 2014 como contraponto. (f) O Sul repensa: Quijano 2000; Giorgini e Iumatti 2025, p. 5 (da dependência ao decolonial); Antunes de Oliveira e Kvangraven 2023 (Dakar, Samir Amin). (g) China: Arrighi 2008 (de fora); autores chineses fora do Zotero: Lin Chun 2013, Wen Tiejun 2021, Cheng Enfu e Lu Baolin 2021 (posição oficial: marcar como tal); Hung Ho-fung 2015 como ponte.}

Institucionalismo e dependência (as instituições chave)

\todoaiflag{}{BIBLIOGRAFIA (IA, 2026-10-29, pasta A.3) para "Institucionalismo e dependência". (a) Gordon 1982 (Journal of Economic Issues) é o único texto que trata literalmente do tema: lê o "tripé" de Cardoso como instituições (p. 569), mas critica a dependência pelo institucionalismo vebleniano (pp. 571-573). Usar como contraponto. (b) Instituições-chave a nomear: Estado (Evans 1995; Rueschemeyer e Evans 1985); alianças de classe e capital estrangeiro (Evans 1979; Cardoso e Faletto 1970); regime de trabalho (Silver 2003). (c) Fio condutor polanyi para evans: Evans 2023, pp. 4-5. (d) Lacuna do esboço: dizer que tipo de institucionalismo a dependência é. Sugestão: institucionalismo histórico-estrutural, a "nova síntese da economia política histórico-estrutural" (Bastos 2025, pp. 3 e 15, a partir de Bambirra), o que costura este bloco com Polanyi e Evans acima. (e) Os Tópicos 1 a 7 (IA, 2026-10-18) se sobrepõem a este esboço (Tópico 2 é a dependência; 3 e 4, desenvolvimento e classes; 5, Polanyi e Evans): sugiro fundi-los e apagar os redundantes.}

\section{Teoria Sociológica sobre Desindustrialização}\label{teoria-socioluxf3gica-sobre-desindustrializauxe7uxe3o}

\subsection{Desindustrialização no Mundo}\label{desindustrializauxe7uxe3o-no-mundo}
NOTA\footnote{Esse capítulo utiliza trechos levemente alterados de artigos publicados 
anteriormente durante a pesquisa de doutorado, a saber "Desafios da reindustrialização brasileira",
publicado como nota técnica pelo Boletim da USCS e a revisão de literatura "Desindustrialização: 
balanço da literatura recente na sociologia e história", a ser publicado na Revista BIB ANPOCS}
\subsubsection{Introdução}

Esse capítulo é uma revisão de literatura sobre processos de
desindustrialização. Seu objetivo é caracterizá-los em aspectos
sociológicos e históricos, identificando tanto suas características
generalizáveis como suas especificidades temporais e geográficas. Visa
também atualizar o leitor sobre os principais desenvolvimentos da
literatura de língua inglesa recente sobre o tema nas áreas de
sociologia e história.

O termo ``desindustrialização'' é comumente tratado como uma redução da
participação da manufatura e da indústria geral na economia
\citep[p.~1]{turner_industrialization_2017} ou
uma redução do emprego industrial em proporção ao estoque de empregos
(\citeauthor{ritzer_deindustrialization_2012}, \citeyear{ritzer_deindustrialization_2012}, p.~1; \citeauthor{ritzer_deindustrialization_2020}, \citeyear{ritzer_deindustrialization_2020}, p.~1; \citeauthor{orum_deindustrialization_2019}, \citeyear{orum_deindustrialization_2019}, p.~1). Em geral é
identificada pela mudança em indicadores de participação do setor
industrial no PIB
\citep{rowthorn_growth_1999} conjugadamente com indicadores de emprego no nível nacional
\citep{tregenna_characterising_2009}.
Também pode ser abordado a partir de suas manifestações locais, como o
fechamento de fábricas
\citep{clarke_closing_2015,hodson_titanic_2019}, regionais e setoriais
\citep{berger_burdens_2017,durazzi_social_2020}. Há na literatura também uma
crescente defesa de uma visão menos economicista do fenômeno, seja
abordando seus múltiplos significados para trabalhadores e comunidades
afetadas \citep{cowie_beyond_2003} ou tentando situá-lo como um processo histórico de
encerramento da sociedade industrial
\citep{strangleman_deindustrialisation_2017}.

Neste trabalho propõe-se como definição de desindustrialização \textbf{a
transformação social associada à redução relativa da participação da
indústria na capacidade produtiva e no estoque de empregos em dado
território, que passa a se dar de forma acelerada e recorrente a partir
da segunda metade do século XX.} Essa mudança tem dimensões econômica,
política e cultural e é mediada por instituições tempo-espacialmente
enraizadas. Essa definição é resultado da revisão de literatura
apresentada, mas também da proposição de
\citep{santos_desenvolvimento_2016} de
abordar o desenvolvimento econômico enquanto mudança social numa
perspectiva socioantropológica.

Enfatiza-se aqui, entretanto, seu caráter de transformação social
historicamente situada, como defendido em \citet{strangleman_deindustrialisation_2017} e \citet{lawson_making_2020}, e a
necessidade de conectar os estudos sobre trabalhadores e suas
comunidades às transformações de nível macro na economia política e na
divisão internacional do trabalho, como propõe High
\aiflag{(High, 2023, p. 56)}.
Destaca-se a importância de compreender a aceleração dessa transformação
social, resultante da interação complexa de múltiplos fatores
contextuais, em vez de buscar causas únicas.

Entre as razões para investigar o fenômeno estão seus efeitos sobre e
relações com a política
\citep{iversen_causes_2000}, comércio
\citep{alderson_explaining_1999,alderson_globalization_2002}, emprego
\citep{beatty_long_2020,tomlinson_-industrialization_2020}, meio ambiente
\citep{feltrin_noxious_2022} e seus impactos nas identidades dos indivíduos e comunidades
afetados por ele
\citep{linkon_half-life_2018,mah_industrial_2012,nayak_displaced_2006,perchard_broken_2013}. Inicialmente identificada com o
rust belt americano e regiões europeias como o Ruhr, Alemanha, ou norte
da Inglaterra \citep{emery_geographies_2019,high_wounds_2013,strangleman_introduction_2013}, a discussão tem
ganhado crescente espaço fora do Norte global, sendo seu aparecimento
``prematuro'' visto como um desafio para países de industrialização
tardia \citep{rodrik_premature_2016}.

No Brasil, as discussões centram-se sobre sua ocorrência, causas e
consequências \aiflag{(Morceiro e Guilhoto, 2020)}, mas o debate acaba centrado no campo da
economia, limitando a compreensão da complexidade do fenômeno. Como
exceção, Santos et al
\aiflag{(Santos et al., 2024)} apontam sua relação com mudanças na dinâmica eleitoral
brasileira. Por sua vez, políticas públicas como a ``Nova Indústria
Brasil'' partem da premissa da existência da desindustrialização e visam
revertê-la ou amenizá-la como parte de uma estratégia de desenvolvimento
econômico \citep[p.~5]{cndi_nova_2024}

O texto a seguir está dividido em quatro partes que abordam: 1 -
critérios de seleção e características da revisão de literatura; 2 -
características econômicas, culturais e políticas da
desindustrialização; 3- a temporalidade, lugares e grupos sociais
privilegiados no seu debate; 4 - suas causas, consequências e relações
territoriais e setoriais; 5- discussão e conclusões.

\subsubsection{Características da revisão de literatura}

A revisão de literatura realizada é exploratória e crítica
\citep[p.~15]{jesson_doing_2011}, incorporando elementos de uma revisão sistemática
\citep[p.~25]{efron_writing_2019}, mas com flexibilização da cobertura do material. Foram
utilizados termos de busca na \emph{Web of
Science}\footnote{Os termos de busca também são denominados comumente
  como \emph{querys} e tipicamente permitem o encadeamento booleano de
  combinações de termos que devem estar presentes ou serem excluídos dos
  resultados.}\footnote{A \emph{query} utilizada teve o seguinte
  formato: \emph{(deindustrialization OR deindustrialisation OR
  de-industrialization OR de-industrialisation OR reindustrialization OR
  reindustrialisation OR re-industrialization OR re-industrialisation OR
  ``manufacturing decline'' OR ``decline in manufacturing'' OR
  ``industrial closings'' OR ``shift to service economy'' OR
  ``transition to service economy'')}, sendo os termos procurados no
  título, resumo ou palavras chave dos artigos (critério ``topics'' pela
  nomenclatura da \emph{Web of Science}).} para encontrar 205 artigos em
inglês nas áreas de sociologia e história, publicados entre 2014 e 2023
e revisados por pares.

Esses artigos foram ordenados por número de citações e avaliados quanto
ao seu tópico central, abordagem do conceito de desindustrialização,
dimensões consideradas e metodologia (cf. tabela 11 abaixo). Vinte e
três artigos foram selecionados com base nesses critérios. Em seguida,
foram extraídas as referências citadas pelo conjunto de 205 artigos,
resultando em uma lista das fontes bibliográficas mais utilizadas nos
últimos dez anos. Foram selecionados os artigos revisados por pares mais
relevantes, adicionando assim vinte artigos e os cinco livros mais
citados ao \emph{corpus}. No total, tem-se a seleção sistemática de 43
artigos e cinco livros como base para a revisão de literatura. Por fim,
critérios não sistemáticos foram utilizados para preencher lacunas
específicas relativas ao caso brasileiro e contextualizar a discussão na
introdução acima.

\begin{table}[htbp]
\centering
\caption{Critérios de priorização de artigos}
\label{tab:cap02-criterios-priorizacao}
\csvautotabular{tables/cap02/criterios-priorizacao-artigos.csv}
\end{table}

Nota: (+) indica aspecto positivo; (-) indica aspecto negativo; pontuação de 0, 0.5 ou 1 em cada critério; selecionados artigos que pontuaram 3 ou mais.

Fonte: autoria própria

Realizando uma inferência a partir da revista de publicação, ou da
formação acadêmica do autor, pode-se agrupar os artigos da seguinte
maneira: sociologia (15), história (18), economia (4), ciência política
(3), estudos culturais (1), \emph{social policy} (1) e geografia (1). Os
livros, a partir da formação dos autores, podem ser classificados em:
sociologia (1), história (2), economia (1) e estudos de literatura (1).
Tal agrupamento, entretanto, não dá conta completa da natureza
interdisciplinar dos trabalhos do campo. Tal interdisciplinaridade se
reflete na heterogeneidade das linhas teóricas utilizadas, estando
presentes abordagens neoclássicas
\citep{scheiring_deindustrialization_2023}, marxistas
\citep[p.~997]{high_wounds_2013},
institucionalistas
\citep{lee_institutional_2016}, da
história econômica
\citep{van_nederveen_meerkerk_challenging_2017} e da sociologia da estratificação e mobilidade social
\citep{kollmeyer_trade_2018},
entre outras.

Do ponto de vista metodológico, predominam na literatura trabalhos
primordialmente qualitativos. Há forte presença de trabalhos que
utilizam metodologias de história oral
\citep{clark_voices_2020,clarke_closing_2015,hodson_titanic_2019}, entrevistas
\citep{hoekstra_iconic_2020,nettleingham_beyond_2019,scheiring_left_2020,scheiring_deindustrialization_2023} e em
menor grau análises documentais
\citep{phillips_closure_2015,phillips_deindustrialization_2019}, etnografias
\citep{lewis_regenerating_2016} e
análises institucionalistas comparativas
\citep{lee_institutional_2016}.
Alguns trabalhos empregam métodos mistos, combinando análises históricas
com análises de dados quantitativos
\citep{tomlinson_-industrialization_2020,van_nederveen_meerkerk_challenging_2017}.

Já trabalhos de ênfase primordialmente quantitativa são minoria, se
destacando preocupações com o impacto econômico da desindustrialização
na desigualdade econômica e mobilidade social
\citep{berger_industrial_2022,kollmeyer_trade_2018} ou suas causas do ponto de vista econômico
\citep{kollmeyer_trade_2018,rodrik_premature_2016}. Estes são na maioria baseados em técnicas de
regressão multivariada. Por fim, alguns trabalhos têm preocupação
primordialmente teórica
\citep{strangleman_deindustrialisation_2017,tomlinson_-industrialization_2016} enquanto outros são revisões de literatura
\citep{emery_geographies_2019,high_wounds_2013,lawson_making_2020,strangleman_new_2014,strangleman_introduction_2013}.

\subsubsection{Características econômicas, culturais e políticas do processo de desindustrialização}

Diversas maneiras pelas quais o processo de desindustrialização é
compreendido na literatura podem ser agregadas a partir da ênfase em
seus aspectos econômicos, culturais ou políticos. Obviamente a maioria
dos trabalhos não adere estritamente a uma das categorias, em geral
sendo valorizada a interação de duas ou mais dessas categorias.

No que concerne à dimensão econômica, a redução do emprego industrial é
o principal indicador do fenômeno na literatura examinada
(\citeauthor{beatty_long_2020}, \citeyear{beatty_long_2020}; \citeauthor{berger_industrial_2022}, \citeyear{berger_industrial_2022}; \citeauthor{clark_voices_2020}, \citeyear{clark_voices_2020}, p.~42; \citeauthor{phillips_deindustrialization_2019}, \citeyear{phillips_deindustrialization_2019}; \citeauthor{scheiring_left_2020}, \citeyear{scheiring_left_2020}; \citeauthor{tomlinson_-industrialization_2016}, \citeyear{tomlinson_-industrialization_2016}, p.~84; \citeauthor{van_nederveen_meerkerk_challenging_2017}, \citeyear{van_nederveen_meerkerk_challenging_2017}), em diversos casos analisada a partir da
experiência de fechamento de fábricas
\citep{clarke_closing_2015,hodson_titanic_2019,lewis_regenerating_2016}, saída de empresas de regiões
\citep{highsmith_beyond_2014} ou
declínio de setores industriais particulares, como carvão e aço
\citep{berger_burdens_2017}. Alguns autores tratam o processo, de forma menos
específica, como uma redução da capacidade produtiva
\citep{bluestone_deindustrialization_1982}, um encolhimento dos setores industriais centrais
\citep{durazzi_social_2020} ou uma redução da produção -- \emph{output} no original --
industrial \citep[p.~1219--1220]{van_nederveen_meerkerk_challenging_2017}.

Para \citet{bluestone_deindustrialization_1982}, \citet[p.~1]{kollmeyer_trade_2018} e \citet[p.~40]{clark_voices_2020}
a contração do setor industrial está relacionada à crescente mobilidade
de capital e para \citet[p.~199]{tomlinson_-industrialization_2016,tomlinson_-industrialization_2020}
ao crescimento do setor terciário e transição para uma economia de
serviços. Tal transição é associada à crescente terceirização e aumento
do trabalho precário em \citet[p.~79]{lee_institutional_2016} e a
salários mais baixos em \citet[p.~107]{winant_hard_2019}. Ainda
sobre o mercado de trabalho, há ênfase numa crescente dualização, a
separação deste em camadas que mantêm certo grau de proteção social e
uma crescente camada de trabalhadores submetidos a piores condições de
trabalho e menores salários
(\citeauthor{durazzi_social_2020}, \citeyear{durazzi_social_2020}, p.~120; \citeauthor{lee_institutional_2016}, \citeyear{lee_institutional_2016}, p.~86).

A conceituação do fenômeno, porém, vai além, incorporando dimensões
culturais e políticas. É notável a influência programática de Cowie e
Heathcott nos estudos de desindustrialização ao propor repensar sua
``cronologia, memória, relações espaciais, cultura e políticas'' {[}e
analisar{]} ``os múltiplos significados de uma das mais importantes
transformações do século XX''
\citep[p.~2]{cowie_beyond_2003}\footnote{Tradução própria}.

Dentro da dimensão cultural, algumas importantes considerações são
sobre como essa mudança representa uma transição de um modo de vida e
cujos impactos se estendem longamente. \citet{linkon_half-life_2018} vai enfatizar
essa continuidade pela expressão "meia vida da desindustrialização",
destacando os efeitos de longo prazo sobre trabalhadores e comunidades
afetados pelos fechamentos de fábricas. \citet{mah_industrial_2012} irá olhar o
constante declínio dessas comunidades sob a perspectiva de uma
"ruinação", enfatizando assim seu aspecto de processo.

Já \citet{strangleman_deindustrialisation_2017} avança a
discussão propondo a leitura dessa transformação como um momento
histórico simétrico à industrialização. Em sua leitura de Thompson
(1961 \emph{apud} \citeauthor{strangleman_deindustrialisation_2017}, \citeyear{strangleman_deindustrialisation_2017}) ele enfatizará este
momento como uma transição que carrega consigo as marcas do período
anterior, um fechamento da era industrial
\citep[p.~467]{strangleman_deindustrialisation_2017}. De forma similar, \citet[p.~694]{hoekstra_iconic_2020}
encapsula as discussões sobre transformações relativas às comunidades
como constituindo a perda de uma forma de vida. Scheiring e King
\citep[p.~144]{scheiring_deindustrialization_2023}vão propor ver a indústria como uma instituição social
cuja sua dissolução provoca uma desintegração com amplas implicações
para o tecido social.

Memória e nostalgia são também mobilizados
\citep{clark_voices_2020,clarke_closing_2015,lewis_regenerating_2016,nettleingham_beyond_2019}. Porém, longe de
estarem ligados meramente ao passado, esses conceitos ressaltam como os
atores sociais mobilizam as narrativas do passado no presente enquanto
recursos, possibilitando assim a transformação do futuro. Dizem respeito
menos ao passado e mais aos ``presentes alternativos''
\citep[p.~119--123]{clarke_closing_2015} possíveis de existir e aos futuros possíveis de serem
imaginados.

Já sobre a dimensão política, de conflito e coordenação, \citet[p.~551]{phillips_closure_2015}
destaca o caráter desejado e politizado da desindustrialização.
Mostra, pela história do fechamento de uma mina de carvão, como as
decisões -- e propensões -- dos atores políticos e das organizações
econômicas são baseadas em suas relações sociais e tensões e que o
processo sempre pode ser gerido de formas distintas
\citep[p.~571--572]{phillips_closure_2015}. Já Clarke
\citep[p.~108]{clarke_closing_2015} mostra como a desindustrialização está imbricada numa
narrativa de ``modernização'' pela qual o fenômeno é considerado
inevitável, funcionando como estratégia por parte de representações
oficias do poder público de forma a despolitizar as discussões sobre
projetos relativos aos antigos territórios industriais.

Por sua vez, \citet[p.~45]{clark_voices_2020}
apontam os desequilíbrios de poder e distintas experiências de
atores da sociedade civil, econômicos e políticos em regiões passando
por entradas ou saídas de investimentos. Já numa dimensão mais ampla,
\citet{tomlinson_-industrialization_2016}
aborda a desindustrialização como uma metanarrativa explicativa das
transformações da Grã-Bretanha no século XX e destaca suas consequências
para a economia política e bem-estar econômico
\citep[p.~86]{tomlinson_-industrialization_2016}. Como um todo, há assim uma noção de desindustrialização como
uma disputa pertencente à economia política na qual claramente alguns
atores -- do mercado e Estado -- possuem posição privilegiada em relação
aos trabalhadores e suas comunidades.

Por fim, a gestão da desindustrialização no tempo presente é
tratada mais centralmente em dois conceitos político-culturais. \citet[p.~64]{rhodes_youngstowns_2013} e \citet{byrne_industrial_2002} mobilizam a
ideia de estrutura de sentimentos, de Raymond Williams, como forma de
descrição da maneira como o processo histórico é percebido no dia a dia
por parte daqueles envolvidos. Já a expressão ``economia moral'', em
\citet{gibbs_moral_2018} e
\citet{tomlinson_re-inventing_2011}, em
grande parte referenciando \citet{thompson_making_1966}, diz respeito
a aspectos normativos, difusos nas comunidades, sobre como a gestão da
transformação econômica deveria se dar, especialmente do ponto de vista
da ação estatal e da mitigação, ou não, dos impactos sobre trabalhadores
e comunidades.

\subsubsection{Lugares, periodização e personagens}

A maior parte das análises concentra-se em países com alto grau de
renda \emph{per capita} e industrialização antiga, em especial na Europa
e América do Norte. As mais frequentes remetem à Grã-Bretanha, aparente
reflexo de sua importância histórica no processo de industrialização.
\citet{tomlinson_-industrialization_2016}
toma como recorte a Grã-Bretanha como um todo a partir de 1945,
\citet{nettleingham_beyond_2019}
dá ênfase à cidade inglesa de Faversham, enfatizando os precedentes
históricos dos fechamentos em massa dos anos 80 e \citet{beatty_long_2020} expandem o
escopo para um conjunto de antigas cidades industriais inglesas. Já
\citet{lewis_regenerating_2016}
trata em específico da região leste de Manchester. Aparecem ainda
discussões sobre a Irlanda do Norte
\citep{hodson_titanic_2019} e
Escócia \citep{perchard_broken_2013,phillips_deindustrialization_2013,phillips_closure_2015,phillips_deindustrialization_2019}.
Há, ainda, debates centrados na Europa continental, como na região do
Ruhr, na Alemanha
\citep{berger_burdens_2017}, França
\citep{clarke_closing_2015} e
Hungria pós transição para uma economia de mercado
\citep{scheiring_left_2020,scheiring_deindustrialization_2023}.

Também central é o caso dos EUA, tratado na escala nacional
\citep{berger_industrial_2022,kollmeyer_trade_2018}, regional, especialmente quanto ao que se
convencionou chamar de ``\emph{rust belt''} americano
\citep{bluestone_deindustrialization_1982,high_industrial_2003} e local, com particular atenção a cidades como
Flint \citep{highsmith_beyond_2014} e Pittsburgh
\citep{winant_hard_2019}.
Ainda na América do Norte, \citet{clark_voices_2020} enfocam a
Nova Escócia, Canadá, ecoando a presença do país como um dos pioneiros
nos estudos de desindustrialização
\citep[p.~994]{high_wounds_2013}.

Alguns trabalhos destacam-se pelo caráter comparativo internacional,
como: \citet{mah_industrial_2012}
que traz comparações de cidades americana, inglesa e russa; \citet{durazzi_social_2020},
centrando a discussão nos contextos Austríaco e Alemão; e \citet{lee_institutional_2016}, que
analisa comparativamente Japão, Coreia do Sul e Taiwan. Este último se
mostra uma das poucas exceções a fugir do eixo Norte Atlântico.

Do ponto de vista da escala analítica, a maior parte dos artigos
trabalha com a escala do Estado-Nação ou país. Também são frequentes
artigos tomando como escala analítica determinada cidade em específico,
às vezes expandindo sua análise para a região metropolitana, como em
\citet{highsmith_beyond_2014}. Em alguns
exemplos mais raros há ênfase regional, como em \citet{berger_burdens_2017}
sobre a região do Ruhr ou \citet{scheiring_left_2020}, o qual
seleciona um conjunto de cidades húngaras tipicamente industriais antes
da transição para economia de mercado. Alguns trabalhos deslocam a
escala ainda mais para a unidade industrial, como \citet{hodson_titanic_2019}, no qual
há ênfase no encerramento de um icônico estaleiro e \citet{phillips_closure_2015}, onde se
descreve em detalhe o fechamento de uma mina de carvão, ambos na
Escócia. Dada a natureza multiescalar dos processos econômicos, uma
lacuna dessa literatura reside na falta de trabalhos operando em
diversas escalas analíticas simultaneamente ou com abordagens de rede.

É Importante salientar como a escolha da escala analítica muitas
vezes é determinante na identificação da existência de processos de
desindustrialização, em especial em situações que envolvem realocação
produtiva. Um dos melhores exemplos é como \citet{bluestone_deindustrialization_1982} enfatizam
simultaneamente a ocorrência de desindustrialização no norte dos EUA
enquanto observam a industrialização no sul do país. Também pioneiro, e
poucas vezes repetido, é o modo como esses mesmos autores abordam os
impactos sociais dos dois processos, mostrando-os como duas faces de um
mesmo fenômeno relacionado à crescente mobilidade de capital.

Essa transformação, entretanto, não segue um ritmo linear ou
simultâneo em todos os territórios. Considerando a escala nacional para
fins didáticos, as análises do caso americano apontam um intenso período
de desinvestimento produtivo nos anos 60 até os anos 90
\citep{bluestone_deindustrialization_1982}, com intensificação nos anos 70
(\citeauthor{highsmith_beyond_2014}, \citeyear{highsmith_beyond_2014}, p.~42; \citeauthor{strangleman_deindustrialisation_2017}, \citeyear{strangleman_deindustrialisation_2017}, p.~474; \citeauthor{winant_hard_2019}, \citeyear{winant_hard_2019}). \citet{kollmeyer_explaining_2009} chega a
pontuar seus sinais já no pós Segunda Guerra Mundial, enquanto \citet{highsmith_beyond_2014} situa-a como
centrada nos anos 1960 em diante, porém ligado intimamente a
descentralização do investimento produtivo já a partir dos anos 1940.
Sobre o Canadá, \citet{clark_voices_2020} destacam
um processo a partir dos anos 1970 em diante, intensificado nos anos
1980.

Os casos da Grã-Bretanha e europeus continentais exibem padrão
similar. No primeiro, há evidências iniciais de desindustrialização
cujas causas remontam ao período prévio a 1945
\citep{tomlinson_-industrialization_2020}, mas cujos desenvolvimentos passam a ser claramente vistos em
fins dos anos 1950, abrangendo a Inglaterra, a Escócia e a Irlanda do
Norte e se aceleram ao longo dos anos 1960 e 1970
\citep{hodson_titanic_2019,phillips_closure_2015,tomlinson_-industrialization_2016}. O período de 1979 a 1982 é,
particularmente, crítico para a Grã-Bretanha como aponta \citet[p.~87]{tomlinson_-industrialization_2016}, de
1975 a 1981 para a Escócia
\citep{phillips_deindustrialization_2019}, se mantendo intenso ao longo dos anos 1980
\citep[p.~2]{beatty_long_2020}.

Tanto EUA como Inglaterra, demonstram, então, alguns indícios de uma
intensificação da desindustrialização e seus impactos com a eleição dos
governos Reagan (1980) e Tatcher (1979) respectivamente, mas há uma
avaliação de que está é prévia a esses governos e já vinha se acelerando
mesmo antes da ascensão de governos classificados como neoliberais
\citep{tomlinson_-industrialization_2016}. Winant argumenta inclusive sobre como os incrementos de
custos relativos à seguridade social decorrentes da desindustrialização
foram na verdade elemento importante nas disputas da economia política
que levaram à ascensão de tais governos
\citep[p.~113--114]{winant_hard_2019}.

Indo em direção a Europa continental, na Alemanha duas ondas se
sobrepõem: uma ligada à indústria de extração de carvão nos anos 1950 em
diante e outra à siderurgia, dos anos 1970 em diante na região do Ruhr
\citep{berger_burdens_2017}. No entanto, o território exibiu maior gradualidade
no fechamento de fábricas, prolongando o processo e mitigando seus
custos sociais
\citep{berger_burdens_2017}. Já na França, é possível ver um sinal de como a
aceleração maior deu-se apenas em fins dos anos 1980 e início dos anos
1990 \citep{clarke_closing_2015}.

O Leste Europeu e o Leste Asiático, por sua vez, apresentam uma
periodização distinta. Na Hungria esse processo se deu de forma
acelerada ao longo dos anos 1990, possivelmente contribuindo para uma
intensificação de seus impactos
\citep{scheiring_left_2020,scheiring_deindustrialization_2023}. Já no Leste Asiático, há identificação de
desindustrialização dos anos 1990 em diante, porém, sua ocorrência é
mais gradual e mitigada pela ação estatal
\citep{lee_institutional_2016}. Há
notável ausência, nessa literatura, de discussões sobre o caso Chinês.

Em resumo, apesar de não significar a inexistência de casos
temporalmente anteriores, como mostram outros trabalhos
\citep{clingingsmith_deindustrialization_2008,pamuk_ottoman_2011}, há um consenso na
literatura sobre sua maior intensidade ocorrendo entre os anos 1960 e
fins dos anos 1980, especialmente na Europa Ocidental e EUA. Territórios
caracterizados por uma industrialização mais tardia\footnote{Desindustrialização
  mais tardia aqui refere-se aos países onde seu processo de
  industrialização ocorreu durante o século XX, em contraposição a
  processos de desindustrialização nos séculos XVIII e XIX
  \citep{mello1982}}
ou com nível de renda \emph{per capita} menor que os casos americano
e europeu parecem ter seus processos situados a partir dos anos 1980,
com intensificação ao longo dos anos 1990 e 2000, dialogando com o
diagnóstico de \citet{rodrik_premature_2016}.

Da mesma forma, algumas categorias de indústria são privilegiadas no
debate, particularmente em função de sua relevância para os territórios
enfatizados. O exemplo alemão, como já apresentado, mostra um processo
de desindustrialização iniciado na indústria extrativa, de
carvão,\footnote{Vale pontuar a relação entre a formação da União
  Europeia e a gestão dos recursos de carvão e aço no continente, sendo
  a primeira resultado direto da formação da Comunidade Europeia do
  Carvão e do Aço (CECA) em 1951
  \citep{valls_european_2016}} que
posteriormente se estende para a indústria de transformação pesada, a
siderurgia \citep{berger_burdens_2017}. A indústria de carvão também é a indústria
central examinada no exemplo Britânico
\citep{tomlinson_-industrialization_2016}. Há
uma lacuna nos estudos a respeito da composição típica da indústria em
cada território e seus impactos. Assim, não pode ser excluída a
possibilidade de fatores de caráter subsetorial influenciarem as
trajetórias observadas.

Já nos EUA e no leste asiático predominam estudos sobre o fechamento de
fábricas e unidades produtivas ligadas à manufatura. No \emph{rust belt}
americano há um aspecto emblemático do fechamento de fábricas ligadas à
indústria automotiva, mas também de serralherias, fábricas de
eletrodomésticos e outros
\citep{cowie_beyond_2003,high_industrial_2003}. Há uma lacuna nos estudos a respeito da composição
típica da indústria em cada território e seus impactos. Assim, não pode
ser excluída a possibilidade de fatores de caráter setorial
influenciarem as trajetórias observadas.

Também há uma notada ausência de discussão a respeito do setor de
construção e do setor ligado ao fornecimento de energia, água, gás e
saneamento -- em geral denominado \emph{utilities}. Um dos problemas
dessa ausência é justamente dificultar comparações setoriais e também
invisibilizar parte significativa da produção industrial.

\begin{table}[htbp]
\centering
\caption{Resumo dos casos de desindustrialização na literatura}
\label{tab:cap02-resumo-casos}
\csvautotabular{tables/cap02/resumo-casos-desindustrializacao.csv}
\end{table}

Fonte: autoria própria

Pode-se, ainda, recortar a literatura a partir de seus enfoques quanto à
classe, a raça e o gênero. Uma clara conclusão é que os estudos são
majoritariamente baseados na perspectiva de trabalhadores e suas
comunidades, em parte pela formação do campo inicialmente calcada no
ativismo contra o fechamento de fábricas
(\citeauthor{high_wounds_2013}, \citeyear{high_wounds_2013}, p.~995--996; \citeauthor{lawson_making_2020}, \citeyear{lawson_making_2020}, p.~2). Em particular, são enfatizados os
trabalhadores industriais afetados diretamente pelo fechamento de
fábricas ou minas, mas em alguns trabalhos grupos sociais urbanos com
interesses antagônicos as comunidades de trabalhadores, como classes
médias urbanas, podem aparecer, em geral nas discussões sobre projetos
relativos à memória ou reestruturação econômica de comunidades
\citep{hodson_titanic_2019,hoekstra_iconic_2020,lewis_regenerating_2016}.

Trabalhos voltados à dimensão racial parecem concentrados no exemplo
americano. Particularmente, preocupam-se com impactos maiores da
desindustrialização sobre as oportunidades futuras e qualidade de vida
da população negra
\citep{berger_industrial_2022}. Nas metrópoles americanas a desindustrialização tem também
ligação com uma gradual gentrificação, especialmente dos subúrbios, onde
frequentemente o fenômeno está associado à composição racial dos
territórios
\citep{highsmith_beyond_2014,winant_hard_2019}. Nas discussões centradas na Europa foi identificada
pouca discussão sobre como a questão racial modifica os impactos da
desindustrialização nos grupos sociais. Isso a despeito da importância
da população imigrante para reestruturação econômica das regiões como o
Ruhr \citep{berger_burdens_2017}.

Já na questão de gênero, alguns trabalhos abordam esta do ponto de vista
dos impactos da desindustrialização a respeito da masculinidade e dos
papéis de gênero
\citep{nayak_displaced_2006}. Um dos
pontos centrais é como as mudanças nos processos de trabalho são vistas
como mudanças de formas de trabalho mais ``masculinas'' para trabalhos
mais ``femininos'' -- particularmente no setor de serviços. Estes se
refletem em requisitos profissionais distintos, frequentemente
estereotipados em termos de gênero, por exemplo a ideia de que
\emph{soft skills}\footnote{Para discussão sobre e caracterização das
  \emph{soft skills}, ver \citet[p.~99--103]{thompson_working_2010}}
seriam mais tipicamente femininas. Há discussões sobre os impactos sobre
as mulheres de trabalhadores e, mesmo que em menor número, sobre
fechamento de fábricas onde a população empregada era majoritariamente
feminina \citep{clarke_closing_2011}. Há pouca discussão sobre as relações entre desindustrialização e
o aumento da participação feminina na força de trabalho.

\subsubsection{Causas, impactos e instituições mediadoras}

A literatura examinada aponta como os processos de desindustrialização
estão ligados tanto às mudanças no modo de produção industrial que
produzem ganhos de produtividade, como às decisões de desinvestimento e
realocação produtiva das empresas
\citep{kollmeyer_explaining_2009,kollmeyer_trade_2018}. Entre os elementos chave, globalização
\citep{alderson_explaining_1999,alderson_globalization_2002} e automação
\citep{berger_industrial_2022} são mobilizados como causas parciais do fenômeno. No quesito
produtividade, em particular, há uma ênfase sobre como seu crescimento
acima da demanda reduziu a participação da indústria no PIB e estoque de
empregos \citep{rodrik_premature_2016}.

Quando visto desse ângulo, a desindustrialização parece intimamente
ligada à viabilidade econômica da atividade industrial em um dado
território (\citeauthor{hoekstra_iconic_2020}, \citeyear{hoekstra_iconic_2020}, p.~697; \citeauthor{phillips_closure_2015}, \citeyear{phillips_closure_2015}, p.~558). Relaciona-se, assim, tanto à
possibilidade de realizar aquele trabalho com menor número de pessoas
como à facilidade em deslocar a atividade produtiva para territórios de
custos mais baixos, especialmente países mais pobres que observaram
ganhos de produtividade consideráveis
\citep{tomlinson_-industrialization_2016}.
Isso implica numa relação com mudanças tecnológicas relativas à produção
e ao controle, inclusive financeiro. Essas mudanças têm base em avanços
em automação \citep{berger_industrial_2022} e em tecnologias de comunicação.

Longe de um determinismo tecnológico, entretanto, essas mudanças estão
ancoradas em decisões conscientes. O uso particular dado a essas
tecnologias está associado a um conjunto de estratégias gerenciais
voltadas à redução de custos, obtenção de benefícios fiscais e ambientes
regulatórios mais flexíveis, crescente financeirização e redução do
poder de barganha coletiva dos trabalhadores. Essas mudanças gerenciais
se configuram em estratégias de realocação produtiva, terceirização de
atividades e adoção de modelos de trabalho novos
\citep{kollmeyer_explaining_2009,kollmeyer_trade_2018}, como métodos \emph{lean}, produção \emph{just in time} e outras
formas de incrementar flexibilização do trabalho e produção.

Portanto, para além de uma discussão sobre redução de custos ou
produtividade, há uma discussão sobre a disputa pela captura de
valor\footnote{Para discussões sobre estratégia de captura de valor ver
  \citet{henderson_redes_2011}} produzido
advindo desses ganhos de produtividade. \citet{bluestone_deindustrialization_1982} associam esse
movimento a uma resposta corporativa à redução de suas taxas de lucro,
mas seria mais claro colocá-lo como uma resposta corporativa à captura
de valor por parte do trabalho\footnote{Para discussão sobre
  transformação nas estratégias de acumulação e o conflito entre capital
  e trabalhadores ver \citet{harvey_condition_1992}}. Os efeitos
dessa disputa sobre o Estado seriam sentidos futuramente com a aplicação
de políticas neoliberais
\citep{lawson_making_2020,scheiring_deindustrialization_2023}, mas a etapa inicial da disputa pela
apropriação do crescente valor produzido se deu na esfera institucional
dos mercados de trabalho.

É em função dessa dimensão política que o processo não pode ser visto
deslocado da trajetória histórica de cada território. Como enfatizado em
\citet{bluestone_deindustrialization_1982}, mas também
presente em \citet{nettleingham_beyond_2019}, \citet{hodson_titanic_2019} e \citet{clarke_closing_2015}, decisões
conscientes de desinvestimento e desregulação foram gradativamente
responsáveis pela perda de competitividade das indústrias nos
territórios desindustrializados. Portanto, para além de mera
racionalidade econômica, há uma necessidade de compreender essas
mudanças como uma transformação desejada advinda da disputa política.

Em resumo, para além da transformação tecnológica e da oportunidade de
negócios, houve intensa articulação política de forma a tornar a
desindustrialização viável e vantajosa. Isso permitiu uma crescente
captura do valor produzido nos territórios por uma elite econômica, em
detrimento da perda de poder econômico e político do trabalho
organizado, assim como de seus efeitos sobre as comunidades locais.

Essa análise afasta a possibilidade de uma leitura excessivamente
determinista do fenômeno, atribuindo-o como quase consequência natural
do crescimento da renda \emph{per capita}. Essa perspectiva foi típica
da economia neoclássica por muito tempo e inicialmente foi importada
para o repertório sociológico
(Clark, 1957 e Bell, 1973, \emph{apud} \citeauthor{kollmeyer_explaining_2009}, \citeyear{kollmeyer_explaining_2009}, p.~1655--1645). A própria identificação
da ocorrência do fenômeno em graus de renda per capita muito distintos
ao longo do tempo
\citep{rodrik_premature_2016}
aponta a existência de determinantes tecnológicos, organizacionais e
políticos para além da mera ``racionalidade econômica''. Prefere-se,
então, ler a correlação dos processos de desindustrialização com níveis
de renda \emph{per capita} como uma contingência histórica naquele
conjunto de economias nas quais essa mudança social se deu primeiro.

Além de uma diversidade de causas, há uma diversidade de impactos. Como
já dito, o ritmo e temporalidade do processo parecem importar na
dimensão e característica de suas consequências sociais. Pode-se
destacar os impactos da desindustrialização no volume de empregos e, por
consequência, na renda e grau de proteção social desses empregos
\citep{beatty_long_2020}, particularmente na medida em que postos de trabalho criados no
setor de serviço tendem a enfatizar uma dualização do mercado de
trabalho \citep[p.~84--86]{lee_institutional_2016}.

Assim, é central o impacto da desindustrialização no crescimento da
desigualdade, especialmente em conjunto com o enfraquecimento de
sindicatos e outras organizações de trabalhadores
\citep{kollmeyer_trade_2018} e
na piora de diversos indicadores relativos ao bem-estar da população, a
exemplo da ocorrência de depressão, do suicídio, do uso de drogas e
mesmo do aumento de doenças cardíacas
\citep{scheiring_deindustrialization_2023,winant_hard_2019}. Ademais, a perda de empregos tende a afetar
permanentemente a renda e a chance de ocupação de melhores postos de
trabalho pelas populações impactadas
\citep{berger_industrial_2022}, sendo os efeitos na renda mais fortes quando se observa a troca
de ocupação em indústrias tipicamente mais bem remuneradas para outras
de remuneração média inferior
\citep{bluestone_deindustrialization_1982,clark_voices_2020}.

Há, ainda, já citadas consequências relativas à perda de identidades e
modos de vida \citep{beatty_long_2020,linkon_half-life_2018,mah_industrial_2012}. Outro resultado é a
gentrificação e o apagamento da memória social, relacionados aos
projetos de reorientação dos territórios desindustrializados
\citep{lewis_regenerating_2016}. Muitos
trabalhadores sentem a perda de conexão com suas comunidades e senso de
pertencimento aos espaços públicos dos quais historicamente eles foram
parte \citep{clarke_closing_2015,phillips_closure_2015}.

Quanto ao Estado, são significativos os efeitos da desindustrialização
sobre a sua capacidade de arrecadação
\citep{highsmith_beyond_2014} e
no incremento de gastos
\citep{beatty_long_2020}. Em particular, projetos de austeridade fiscal são comuns como
resposta à pressão econômica e social criada pela desindustrialização
\citep{lawson_making_2020}. Para
além, há impactos sobre a percepção de populações sobre grupos políticos
que podem ser duradouras, em geral sendo enfatizada a capacidade do
grupo político no poder de responder ao fenômeno
\citep{phillips_deindustrialization_2019,scheiring_left_2020}. Assim, estes processos parecem
interagir fortemente com transições políticas, seja como causa ou
consequência, e afetam por décadas as forças políticas organizadas em
torno dos territórios desindustrializados.

É possível, ainda, observar como a literatura trabalha a dimensão
relacional, que assume duas grandes ênfases: enquanto alguns trabalhos
destacam o processo como uma relação direcionada da atividade econômica
do setor secundário para o terciário
\citep{beatty_long_2020,clark_voices_2020,clarke_closing_2015,kollmeyer_trade_2018,phillips_deindustrialization_2019,tomlinson_-industrialization_2016,tomlinson_-industrialization_2020,winant_hard_2019}; outros
a enfatizam como uma relação direcionada entre um território dito
``desindustrializado'' e novos territórios que passam por
industrialização
\citep{clarke_closing_2015,hodson_titanic_2019,kollmeyer_trade_2018,strangleman_deindustrialisation_2017,tomlinson_-industrialization_2020}. Nesse
cenário é comum um influxo de imigrantes dos antigos territórios aos
novos \citep{scheiring_deindustrialization_2023,winant_hard_2019}, por vezes em modo pendular
\citep[p.~43]{clark_voices_2020}.

Há na literatura pouco debate sobre as decisões e estratégias dos
territórios para onde se deslocam as atividades industriais, em geral
tendo sua atuação estratégica reduzida à ideia de oferecer mercados de
trabalho com salários mais baixos ou incentivos de ordem fiscal e
regulatória
\citep[p.~474]{strangleman_deindustrialisation_2017}. Isso destoa de achados da sociologia econômica sobre como a
construção social de novos territórios é mais complexa e planejada do
que aparenta, exigindo arranjos sociopolíticos particulares
\citep{rodrigues_o_2007}.

Essa baixa complexidade parece se refletir numa leitura relativamente
rasa da relação entre instituições e desindustrialização. Estados
nacionais são considerados importantes agentes, impactando ritmo e
efeitos sociais, especialmente via uso de mecanismos de seguridade
social \citep{tomlinson_-industrialization_2016}, proteção do emprego
\citep{phillips_closure_2015,phillips_deindustrialization_2019} e via liderança em projetos de
reestruturação econômica
\citep{berger_burdens_2017}. Entretanto, na maior parte dos trabalhos são
considerados monolíticos e não como conjunto de organizações ou
instituições com conflitos internos. Atores ligados ao Estado aparecem
apenas enquanto governos alternantes no nível do executivo federal ou
governos locais enquanto articuladores dos projetos de reestruturação
econômica, mas com pouca atenção ao papel das burocracias estatais,
empresas estatais, legislativo ou judiciário.

Sobre a sociedade civil destaca-se a participação e influência de
organizações ligadas à ação coletiva de trabalhadores, particularmente
sindicatos \citep{hodson_titanic_2019,kollmeyer_trade_2018,phillips_closure_2015,phillips_deindustrialization_2019}, vistos como capazes de, em parte, frear, estabelecer condições
melhores ou reduzir os impactos sobre as populações locais da perda de
empregos industriais. Outro conjunto de instituições relevantes são
organizações de moradores de certos territórios
\citep{nettleingham_beyond_2019},
que são particularmente atuantes na gestão do território
desindustrializado e que buscam influenciar as estratégias voltadas a
reorientar o papel econômico dos territórios
\citep{hodson_titanic_2019,lewis_regenerating_2016}.

Por fim, instituições de caráter econômico, em especial as empresas, são
colocadas como os elementos mais relevantes para o processo de
desindustrialização, sendo ele resultante primariamente de suas decisões
\citep{clarke_closing_2015,kollmeyer_trade_2018}. Entretanto, apesar do enorme peso associado às suas
decisões, análises compreensivas sobre elas são quase ausentes na
literatura. Toma-se como premissa a ação racional e orientada a
interesses de acumulação por parte de empresas cujas decisões seriam
quase ``inevitáveis'' nessa perspectiva. Entretanto, algumas evidências
sugerem que essa relação não é tão simples. Alguns artigos apontam como
a relação entre empresas e territórios é mais complexa
\citep{highsmith_beyond_2014} e
que a decisão de desinvestimento está também ligada a disputas nas
economias políticas dos países
\citep{bluestone_deindustrialization_1982}. Assim, parece precoce naturalizar as decisões de desinvestimento
e realocação produtiva como a única resposta racional ao contexto
apresentado. Assim, outra lacuna diz respeito à escassez de análises
acerca das perspectivas de empresas, elites econômicas ou trabalhadores
em atividade gerencial.

\subsubsection{Discussão e conclusões}

As discussões acima exemplificam como a desindustrialização não possui
uma definição consensual. Antes, possui características gerais
associadas às suas manifestações empíricas, mas estas não
necessariamente são amarradas em uma definição teórica coerente. Isso
torna possível autores divergirem até mesmo sobre a existência do
fenômeno em um dado caso.

Se não há consenso claro sobre o conceito, a maior parte da literatura
examinada compartilha de uma preocupação com a discussão dessa mudança
social do ponto de vista dos atores que o vivenciam mais diretamente,
particularmente daqueles entendidos como membros da classe trabalhadora
industrial, suas famílias e comunidades. São, portanto, tratados como
observadores privilegiados.

Entretanto, discussões de nível analítico macro são deixadas em segundo
plano, talvez com exceção dos trabalhos utilizando análise documental ou
estatística. Isso gera uma lacuna da compreensão sociológica a respeito
do envolvimento de Estado e empresas nas transformações. Assume-se muito
facilmente uma perspectiva baseada na escolha racional para os atores
econômicos, a qual \citet{highsmith_beyond_2014} mostra não ser
suficiente para compreender a formação das decisões de relocação
produtiva ou fechamento de fábricas.

Também acaba relegada ao segundo plano a discussão do poder e seu papel
na conformação do processo de desindustrialização, sendo as formas de
resistência dos trabalhadores quase sempre entendidas como insuficientes
para alterar significativamente o curso dessa transformação
\citep{high_fruitless_2022}. Com
poucas exceções, como \citet{berger_burdens_2017}, há pouca
preocupação sobre a inter-relação entre movimentos de trabalhadores e
comunidades e sua capacidade efetiva de influenciar os destinos dos
territórios. Dada, também, a pouca frequência de trabalhos centrados no
Estado ou empresas, o poder acaba sendo uma categoria pouco
problematizada. As transformações impõem-se sobre as comunidades as
quais apenas reagem a elas, sem discussão dos interesses e estratégias
orientadores das decisões dos atores organizacionais envolvidos.

Apontar esse viés não implica defender uma perspectiva elitista ou um
foco excessivo em instituições formais. Antes disso, é justamente a
compreensão tanto do Estado quanto das elites econômicas enquanto grupos
de atores complexos, contraditórios e socialmente enraizados que permite
tornar estes acessíveis à ação coletiva da sociedade civil. Se Estado e
empresas não forem vistas em sua realidade, compreendidos em suas formas
de articulação e uso do poder, nunca será claro as maneiras como as
populações afetadas são capazes de reagir.

Esse quadro dialoga com análises como a de \citet{high_fruitless_2022}, que reflete
sobre a ação de trabalhadores ser muitas vezes ``infrutífera''.
Especialmente diante de uma perspectiva que deseja enfatizar as
contradições dessa transformação com esforços de neo ou
re-industrialização, como no caso brasileiro, é preciso tornar Estado e
empresas analiticamente ``vulneráveis'' aos contextos sociais. Isso
permite, por exemplo, ver como mudanças econômicas se refletem em novos
arranjos e aprendizados institucionais operando nos territórios
\aiflag{(Ramalho, 2005)}.

Carece, portanto, de uma discussão mais aprofundada sobre como a ação
social desses atores coletivos dotados de uma dimensão de poder, Estado
e empresas em particular, é afetada pelas instituições nacionais e
regionais. A ausência dessa discussão é não compreender parte importante
do argumento de \citet[p.~15]{bluestone_deindustrialization_1982} de que
``desindustrialização não apenas acontece''\footnote{Tradução própria}.
Isso implica discutir as estratégias de acumulação das organizações
econômicas e a ação do Estado na forma de suas políticas públicas e
organização burocrática.

Adicionalmente, se a bibliografia analisada permite uma boa leitura para
os casos das nações do Norte Atlântico, não é tão capaz de lidar por
completo com a realidade do fenômeno em territórios com trajetórias
históricas distintas, como o Brasil. Quando o processo de
desindustrialização ocorre simultaneamente com industrialização tardia,
\emph{neo} ou re-industrialização ou em comunidades nas quais nunca a
era industrial não teve o mesmo impacto, as perspectivas teóricas
oriundas de como esse processo é vivido nas nações onde ele primeiro se
evidenciou não apresentam o mesmo potencial explicativo. Dito de outra
maneira, a realidade é mais heterogênea do que uma leitura situada na
Europa Ocidental e América do Norte é capaz de dar conta.

É importante também salientar as lacunas e limitações desse próprio
texto. Primeiramente, é necessário ampliar de maneira sistemática o
escopo da análise para abarcar melhor trabalhos recentes das áreas de
economia, geografia e ciência política. Em segundo lugar, é preciso
utilizar critérios não sistemáticos para incluir na discussão mais
trabalhos a respeito dos casos fora do eixo Norte Atlântico -- com
ênfase no Brasil. Por fim, a falta de discussões mais aprofundadas sobre
o impacto da desindustrialização nas relações de gênero e raça, além de
maior reflexão sobre sua dimensão ambiental, são limitações deste
balanço. Essas lacunas apontam a necessidade de continuar, aprofundar e
complementar as discussões aqui presentes em trabalhos futuros.

\subsection{Desindustrialização no Brasil}\label{desindustrializauxe7uxe3o-no-brasil}

\subsubsection{Introdução: breve contexto da desindustrialização brasileira}

Os processos de desindustrialização se tornaram frequentes a partir da
segunda metade do século XX \citep{emery_geographies_2019,high_wounds_2013,lawson_making_2020,strangleman_introduction_2013}. Estes podem ser mensurados pela
redução da participação do setor industrial, em especial da manufatura,
nos Produtos Internos Brutos (PIBs) dos países e em seu estoque de
empregos \citep{tregenna_characterising_2009}. Apesar dessa queda ter sido particularmente
intensa nas economias da Europa Ocidental e dos EUA, ela se manifestou
também em economias menos diversificadas e de menor PIB \emph{per
capita}, em especial na América Latina e Leste europeu \citep{rodrik_premature_2016}.

Dentre todos esses casos, o brasileiro se destaca como um dos mais
intensos e duradouros. Com breves interrupções nessa trajetória, o
Brasil viu a participação da atividade industrial em seu produto bruto
se reduzir desde o início dos anos 1980 até o presente momento \aiflag{(Feijó e
Araujo, 2023; Morceiro, 2021; Oreiro e Feijó, 2010)}. Apenas o setor de
manufatura regrediu de 27\% em 1985 para 17\% nos dados de 1996 e para
11\% em 2019, quando considerada a série temporal corrigida por Morceiro
(2021).

Entre as causas do fenômeno a nível mundial a principal parece ser o
significativo incremento da produtividade do trabalho industrial, que
reduziu fortemente a demanda de trabalhadores \citep{rowthorn_growth_1999}. Ao mesmo tempo, uma dinâmica de realocação produtiva,
especialmente em direção ao Leste Asiático e parte do Sul Global,
intensificou o fenômeno nas economias de industrialização mais antiga
\citep{kollmeyer_explaining_2009}.

Entretanto, tais fenômenos não podem ser compreendidos fora de seu
contexto histórico. Se nos países supracitados o início desse processo
ocorreu nos anos 1950, tal período para o Brasil ainda era de
industrialização e de predomínio de políticas desenvolvimentistas
ancoradas no papel do Estado \citep{perissinotto_ideias_2023} e no capital
estrangeiro \citep{evans_dependent_1979}. Apenas em meados da década de 1970 a
estagnação do crescimento industrial brasileiro começa a exibir seus
sinais. Esta se dá no contexto da crise da dívida externa brasileira e
do crescimento da inflação \citep{hermannetal2015}. Na medida em que o controle
inflacionário se torna a principal preocupação da política fiscal
nacional, a política industrial passa ao segundo plano \citep{cano_brasil_2017},
fato evidenciado pelo abandono do terceiro Programa Nacional de
Desenvolvimento em fins da ditadura militar \aiflag{(Faria, 2016, p. 184--186)}.

A partir dos anos 1980 o país começa a passar por uma efetiva
desindustrialização, ao mesmo tempo que transaciona do regime militar
para a democracia. Em tal contexto de transformação das instituições
políticas, foi contínua a ênfase no combate à inflação como política
fiscal, especialmente após a eleição de Collor{[}1989-1992{]} em 1989,
quando se acelerou um processo de liberalização da economia e abertura
comercial (de Paula, 2022). Após sucessivos planos econômicos mal
sucedidos, o Plano Real consolidou em 1993 as bases de um novo modelo de
crescimento econômico baseado, entre outros pontos, no \emph{superavit}
fiscal, privatizações, juros altos e na paridade real-dólar, até a
flexibilização do câmbio em 1998 \aiflag{(Giambiagi, 2015; Pires, 2022)}. Modelo
esse que, se por um lado logrou o controle hiperinflacionário, não
conseguiu evitar a contínua desindustrialização. Em concomitância de tal
cenário com as reestruturações no modelo produtivo, a ascensão do
chamado Toyotismo \aiflag{(Ramalho, 2010)}, os anos 1990 marcaram significativa
queda no emprego industrial no Brasil \aiflag{(Ramalho e Rodrigues, 2013)}.

Aqui enfatiza-se, portanto, que não se trata de mero ``problema
econômico''. A desindustrialização brasileira precisa ser vista também,
para sua apropriada análise sociológica, como consequência do
direcionamento do capital político brasileiro desde os anos 1970 para a
estabilização inflacionária. Assim, enfatiza-se a capacidade de
articulação política para projetos de desenvolvimento como variável
importante para entendimento da trajetória da economia brasileira e para
análise de seus desdobramentos futuros.

Essa dinâmica acompanhou dois outros fenômenos agravantes.
Primeiramente, desde ao menos os anos 1990 o Brasil vive uma
desconcentração industrial, exibindo uma saída da atividade industrial
da região sudeste, em especial de São Paulo, para outras regiões
brasileiras \aiflag{(CNI, 2021)}. Tal fenômeno intensifica o impacto do processo
de desindustrialização em antigos territórios industriais, além de
frequentemente estar relacionado às disputas fiscais entre estados
\aiflag{(Arbix, 2000)}, via redução de impostos e oferta de subsídios, e ao
incremento da demanda por infraestrutura urbana \citep{bluestone_deindustrialization_1982}, reduzindo a capacidade de resposta fiscal (i.e. a oferta de
recursos financeiros disponíveis) do Estado ao fenômeno.

Em segundo lugar, o incremento do volume de empregos no chamado ``setor
de serviços'', que acompanhou a desindustrialização nos países de
industrialização avançada \aiflag{(Amsden, 2001)}, exibe características
particulares no Brasil. Isso devido, no caso nacional, a baixa
formalização, baixa remuneração, pouca qualificação e pouca perspectiva
de crescimento de carreiras nos principais ``novos'' setores
empregadores \aiflag{(Bridi, Oliveira e Salas, 2024)}. Com particularidades da
atividade econômica emergente em cada período, dos setores de
telemarketing dos anos 1990 até os entregadores de aplicativo nos dias
atuais , o incremento do uso de tecnologia e as inovações gerenciais não
se refletem num incremento na qualificação do trabalho, tornando o
Brasil espécie de paradigma \emph{``bravermaniano''} \aiflag{(Braverman, 1998)}
da desqualificação do trabalho nos tempos contemporâneos.

Para além disso, o movimento acompanha uma tendência de precarização do
trabalho e dissolução das relações de trabalho formais \aiflag{(Ramalho, 2010)}.
No caso brasileiro esse fenômeno encontra evidência nas sucessivas
transformações institucionais \aiflag{(Thelen, 2003)} que a legislação
trabalhista vem sofrendo \aiflag{(Ramalho, Salles Pereira Dos Santos e Jácome
Rodrigues, 2019)}. Tais processos de mudanças institucionais, assim como
as disputas fiscais mencionadas anteriormente, pressionam fiscalmente o
Estado brasileiro ao reduzir a sua capacidade de arrecadação.

\subsubsection{Desafios da reindustrialização brasileira}

Os dilemas fiscais, seja via aumentos dos gastos com seguridade social
ou seja via redução da arrecadação estatal pela redução do emprego
formal, impostos pelos processos de desindustrialização são um elemento
chave para compreender alguns dos principais desafios aos projetos de
reindustrialização. Como Winant (2019) aponta, essas pressões fiscais
tendem a ser refletidas em projetos de austeridade fiscal que, ao menos
parcialmente, buscam mitigar seu impacto. Assim, políticas de
austeridade não podem ser vistas apenas como resultado de uma concepção
ideológica arbitrariamente restritiva de projetos de desenvolvimento.
Elas precisam ser entendidas também como respostas a transformações
estruturais. Por sua vez, essas respostas são enviesadas pelos
interesses das frações de classe dominantes, obscurecendo respostas
alternativas de longo prazo. Isso não implica, entretanto, que
restrições fiscais a projetos de desenvolvimento, e em especial de
reindustrialização, não se imponham.

Dadas essas constrições, tanto as estruturais quanto as institucionais
relativas aos interesses da classe dominante, fica claro que a
capacidade de ação estatal possui limitações. Tais limitações são
relevantes na medida em que a ação do Estado é elemento necessário para
o desenvolvimento econômico, especialmente em economias de
desenvolvimento ou industrialização tardia e que se inserem numa
economia global diante de competidores já ``avançados'' \aiflag{(Amsden, 2001;
Gerschenkron, 1962)}. Especialmente quando se busca reorientar o perfil
das atividades econômicas de um país para além do que seriam suas
``vantagens competitivas'' presentes \citep{evans_embedded_1995}. Significativa
evidência dessa relação entre Estado e desenvolvimento já foi
apresentada para diversas nações \aiflag{(Evans, Rueschemeyer e Skocpol, 1985;
Rueschemeyer e Evans, 1985)}, inclusive o Brasil \citep{evans_dependent_1979,perissinotto_ideias_2023}.

Para analisar tal dinâmica, primeiro pode-se analiticamente separar a
ação do estado na economia em dois eixos: ações voltadas a acrescer o
ritmo de acumulação e ações de caráter redistributivo \aiflag{(Rueschemeyer e
Evans, 1985)}. Dito de outra maneira, a ação do Estado precisa ser
pensada em sua capacidade de: 1- acelerar o ritmo de crescimento
econômico do país ou de setor econômico estratégico; 2- redistribuir os
ganhos econômicos entre diversos grupos sociais. Tal perspectiva é
vantajosa ao avançar de uma dicotomia de ``menos-mais'' Estado para
efetivamente pensar as formas específicas de intervenção estatal \citep{evans_embedded_1995}.

No caso brasileiro, tal dinâmica parece se refletir num \emph{trade-off}
entre esforços estatais para desenvolvimento, que podemos pensar como
estímulo ao investimento, e para ampliação da capacidade de consumo de
determinadas classes sociais. Não se trata aqui de uma oposição direta
entre consumo e investimento, mas de uma limitação na capacidade de ação
do Estado de promover ambos os objetivos \emph{sem incorrer em déficit}
ou tentando reduzi-lo ao máximo, dadas as consequências do déficit
público no modelo de desenvolvimento e ambiente institucional discutido
anteriormente -- pressão inflacionária e aumento de juros.

Essa análise parte da constatação de que, ao menos para o caso
brasileiro, a maior parte dos economistas de orientação ortodoxa como
também de algumas perspectivas heterodoxas influentes, como a novo
desenvolvimentista, consideram a existência de algum grau de equilíbrio
fiscal necessário ao processo de desenvolvimento \aiflag{(Bresser-Pereira,
2012)}. Assim, mesmo que essa política de superávit possa ser revertida
temporariamente, no que se denominaria uma política ``anticíclica'', a
dimensão fiscal tende a ser um fator limitante na capacidade de ação do
Estado no longo prazo.

A discussão acima serve para estabelecer o ponto de que o Estado não é
totalmente livre para livremente estimular investimento e consumo, sendo
necessário balancear em algum grau os dois objetivos. Como resultado se
encontra um problema político: qual consumo e em qual grau ele será
restringido para viabilizar o aumento da poupança e do investimento
produtivo? Como será direcionado o resultado do investimento da
poupança? Não há espaço aqui para reconstituir como, historicamente,
esse conflito se deu nos diversos modelos de desenvolvimento econômico
da história nacional -- por exemplo, o autoritarismo como forma de
ultrapassar a disputa política e do uso de dívida externa para contornar
a limitação fiscal durante a ditadura militar --, mas ele é tomado como
ponto de partida para expressar a complexidade do problema da
desindustrialização e das tentativas de reindustrialização.

Isso porque as mudanças associadas ao processo de desindustrialização
dificultam a resolução do problema político e econômico, que pudesse
propiciar um processo de reindustrialização. Notadamente, desde a
primeira revolução industrial e de forma acelerada a partir da metade do
século XX, a atividade industrial se tornou significativamente mais
complexa do ponto de vista tecnológico e mais intensa em capital
\citep{kollmeyer_explaining_2009,rowthorn_growth_1999}, mesmo que com parcial
exceção de alguns casos no Leste Asiático. Tal processo é coerente com a
crescente dificuldade que economias têm de realizar o ``catch-up''
tecnológico ao longo do tempo \aiflag{(Amsden, 2001)}. Como já apontava
Gerschenkron (1962), cada ciclo de industrialização de um conjunto de
economias parece requerer maior capital e maior atuação do estado para
sua viabilização. Esse maior custo de capital financeiro também implica
maior necessidade de capital político a ser mobilizado, o que pode ser
desafiador quando as elites possuem interesses diretamente associados a
atividades econômicas dependentes do uso intensivo de trabalho, de
caráter extrativista ou rentista.

Ao mesmo tempo, essas características da atividade industrial propiciam
ao trabalho um grau maior de produtividade, especialmente após a segunda
metade do século XX, que implica menor demanda por trabalhadores
\citep{kollmeyer_explaining_2009,rowthorn_growth_1999}. O resultado é que
processos recentes de industrialização, especialmente em setores com
alta perspectiva de crescimento, implicam na alocação de crescentes
volumes de capitais para decrescente número de empregos. A aliança entre
frações de classe que torna diversos projetos de industrialização
possíveis é fragilizada, especialmente em contextos democráticos. Na
ausência de uma elementos viáveis para mobilizar apoio popular e da
classe trabalhadora para avançar com projetos de reestruturação
econômica, o Estado tende a acomodar interesses mais conservadores de
suas elites, ancorados em vantagens competitivas já existentes.

Importante notar que, no passado nacional, essa ``acomodação'' envolveu
muitas vezes autoritarismo significativo, inclusive perante frações da
própria elite. Por exemplo, nos períodos Varguista e da Ditadura Militar
viabilizaram projetos de industrialização baseados na conciliação dos
interesses das elites \citep{perissinotto_ideias_2023}, inclusive as elites
internacionais no segundo caso \citep{evans_dependent_1979}. Mesmo assim, é importante
notar que em ambos os casos algum grau de apoio popular ou ao menos
apaziguação de demandas urbanas era preenchido pela noção do emprego
industrial. Na história brasileira talvez apenas o breve período dos
anos 1950 até o golpe de 1964 tenha conjugado de forma significativa
esse desenvolvimentismo com democracia \citep{barbosa_developmentalist_2020}. Um projeto de
reindustrialização por vias democráticas, se em parte possibilita maior
grau de promoção do bem-estar social e igualdade, ao mesmo tempo exige
maior presença e participação de trabalhadores, suas famílias e suas
instituições de classe para ser viabilizado.

A esse problema adiciona-se mais uma camada: a dificuldade de
coordenação de projetos de reindustrialização entre as diferentes
regiões do Brasil levando em conta suas particularidades. As já citadas
disputas fiscais entre estados são apenas um exemplo. Para além de
dificuldades de conciliar interesses, cabe ressaltar a dificuldade de
promover projetos de integração. Casos bem sucedidos de políticas
estatais para desenvolvimento econômico em geral contam também com a
presença de uma burocracia estatal extremamente capacitada, em algum
grau autônoma e bem integrada com a sociedade \citep{evans_embedded_1995}. Mesmo que
em alguns casos tal capacidade burocrática exista de forma limitada, ela
não aparenta ser em número suficiente e nem contar com o necessário
amparo político-financeiro e capilaridade para desenvolver e coordenar
projetos de nível local e regional na escala exigida \citep{evans_embedded_1995}.

\subsubsection{Os novos argumentos pró-indústria}

O cerne do argumento acima pode ser assim resumido: a reindustrialização
brasileira não é mero problema técnico-econômico que seja superável com
a correta combinação de taxa de juros, câmbio, subsídios públicos,
privatização, estatização ou qualquer instrumento econômico. Antes, ela
existe como um problema político de estabelecer um projeto nacional de
desenvolvimento ancorado na atividade industrial, que enfrenta
obstáculos por um lado devido às característica das elites nacionais e
por outro pela redução da capacidade da atividade industrial de agregar
apoio de fração significativa da classe trabalhadora devido a seu menor
potencial de geração de empregos.

Esse problema não passou ao largo das discussões sobre a
reindustrialização brasileira, mesmo que por vezes tenha sido feito
apenas de forma implícita. Como evidência, um conjunto de argumentos
``pró-industrialização'' tem se formado como resposta a esse fenômeno e
como tentativa de estruturar um projeto coerente de reindustrialização.
Alguns desses argumentos encontraram uma formalização institucional na
nova política industrial, adotada a partir de 2024, denominada Nova
Indústria Brasil \citep{cndi_nova_2024}. É a partir desse referencial que
discutirá-se em seguida alguns desses ``argumentos'', apontando
entretanto alguns de seus limites. Eles dizem respeito ao projeto de
reindustrialização como base de um crescimento econômico sustentável
gerador de empregos ``de qualidade'', como resposta à crise climática,
como ampliador da capacidade exportadora do país, como melhoria do
acesso da população urbana à infraestrutura e serviços de saúde e como
busca de soberania em setores estratégicos \citep{cndi_nova_2024}. Essa
análise fica clara quando observa-se o peso do ``Plano Mais Produção''
dentro do conjunto de instrumentos da nova política industrial, que
``concentra em seus quatro eixos - Mais Inovação, Mais Verde, Mais
Exportação, Mais Produtividade -- para o período de 2023 a 2026''
{[}...{]} ``R\$ 300 bilhões'' \citep[p.~14]{cndi_nova_2024}. Desde seu
anúncio o valor total para o plano foi incrementado em outros 42,7
bilhões \aiflag{(Taiar, 2024)}.

Mas o que destaca-se aqui é menos o valor alocado ao plano, que ainda se
mostra muito inferior a outras iniciativas de apoio ao crescimento, como
os planos Safra \aiflag{(Agência Brasil, 2024)}, e mais as lógicas expressados no
que diz respeito aos retornos esperados à sociedade em função do
desenvolvimento industrial. Dá-se atenção a tais justificativas, pois
elas embasam a discussão sobre a necessidade de uma política pública
sobre o tema e, portanto, são articuladoras entre a ação do Estado e os
interesses dos diversos grupos sociais representados por este.

É preciso esclarecer que o objetivo deste texto não é uma análise
aprofundada da nova política industrial. Em que pese os limites dos
argumentos apresentados acima, o Nova Indústria Brasil representa um
importante esforço de institucionalização de uma política industrial
ancorada simultaneamente em processos de aprendizado institucional e no
impacto direto na vida da população. O que importa aqui, mais
especificamente, é como os objetivos e estratégias delineados no plano
representam não apenas um diagnóstico dos obstáculos ao desenvolvimento
industrial, mas associam a solução destes obstáculos diretamente com os
benefícios que possam ser compartilhados por diferentes grupos sociais.
É assim, um arranjo político que delineia certo conjunto de interesses
como prioritários e suficientes para justificar o uso de recursos
comuns.

Passando diretamente aos argumentos citados, o primeiro diz respeito à
necessidade de não apenas aumentar a produção e produtividade nacional,
mas de um impacto para a população na forma de ``empregos de qualidade''
\citep{cndi_nova_2024}. Em que pese a importância do incremento da
produtividade, sofisticação tecnológica e capacidade exportadora para o
crescimento do país, a articulação de um projeto em torno de ``empregos
de qualidade'' apresenta limitações. A primeira é a já citada limitação
na capacidade de produção de empregos em novas instalações industriais,
dado o paradigma da produção flexível \aiflag{(Ramalho, 2010)}. Uma segunda
problemática é que, mesmo em atividades mais ``sofisticadas'', a
qualidade do emprego não pode ser entendida como resultado necessário ou
mesmo esperado. Como mostra Ramalho (2006), novas unidades produtivas,
com práticas modernas e flexíveis, muitas vezes combinam essas práticas
com modelos de gestão do trabalho tradicionais e frequentemente
autoritários.

Isso implica que, ao passo que produtividade e sofisticação tecnológicas
são elementos necessários para existência de trabalhos de qualidade, não
são causas suficientes para eles \aiflag{(Mahoney, 2003)}. Exemplo recente, a
implementação da unidade produtiva de veículos elétricos em Camaçari tem
levantado embates diante das propostas salariais abaixo das expectativas
de trabalhadores e sindicatos \aiflag{(Gabriel, 2025)}. Tal discussão não é
surpreendente quando tomada em conta a literatura existente sobre a
relação entre processos de \emph{upgrade} econômico e \emph{upgrade}
social \aiflag{(Barrientos, Gereffi e Rossi, 2011)}. Na ausência de elementos
institucionais que garantam a captura do valor produzido por territórios
e comunidades -- como leis trabalhistas, um papel ativo do Estado ou
fortalecimento de organizações representativas dos trabalhadores --
processos de melhoria da produtividade podem não se refletir em
melhorias das condições sociais. No caso brasileiro parece inclusive
existir um movimento na direção contrária com processos de precarização,
através da flexibilização das leis trabalhistas, de redução da
seguridade social e enfraquecimento das organizações trabalhistas.
Assim, reduz a capacidade de mobilização política do emprego industrial
formal.

Tal discussão não significa apenas que projetos de reindustrialização
podem não vir a oferecer retornos sociais e econômicos esperados, mas
que a própria articulação do emprego industrial como base de apoio a
esses projetos fica fragilizada. Como justificar o significativo emprego
de capital econômico e político na manutenção de projetos de longo prazo
que apresentam retornos em empregos quantitativamente inferiores às
alternativas mais precárias e onde o retorno na qualidade desses
empregos encontra-se cada vez mais incerto?

Uma segunda é a questão ambiental e a proposta de posicionar a
reindustrialização como componente de uma agenda de transição
energética. Tal fator é enfatizado inclusive no lugar privilegiado que o
agronegócio tem na pauta do Nova Indústria Brasil \citep{cndi_nova_2024}.
Entretanto, assim como no caso do trabalho de qualidade, o arcabouço
institucional brasileiro parece direcionado de forma contra produtiva a
esse objetivo, vide recentes flexibilizações de regulações ambientais
\aiflag{(León, 2025)}. Nesse sentido, os esforços de construção de uma economia
sustentável e de longo prazo podem ser parcialmente anulados por uma
regulação ambiental que privilegia modelos econômicos menos sustentáveis
e de curto prazo.

Tanto no caso da busca por ganhos de produtividade industrial como na
incorporação do agronegócio via argumento de sustentabilidade, a
ampliação da capacidade de exportação é dado como objetivo privilegiado
\citep{cndi_nova_2024}. Isso é coerente do ponto de vista econômico, já que
a exportação foi importante ferramenta para desenvolvimento industrial
de países de industrialização tardia \aiflag{(Amsden, 2001)}, representando uma
forma de incremento da capacidade de acumulação, via entrada de
capitais, e também reposicionamento do país nas cadeias globais de valor
\aiflag{(Coe \emph{et al.}, 2004; Henderson \emph{et al.}, 2002)}. Entretanto, a
exportação implica, do ponto de vista da economia política, algumas
consequências.

Primeiramente, a orientação da economia às exportações implica, ao menos
no curto prazo, na contenção de consumo. Ao passo que o déficit em
balança comercial implica um grau de consumo adicional, um financiamento
de certo padrão de vida material, um superávit comercial tipicamente
implica um menor consumo doméstico. Isso fica notável ao observar
recentes impactos inflacionários da exportação de produtos alimentícios
\aiflag{(Moura, 2025)} ou na perda de mercado de indústrias locais frente ao
acesso a produtos externos de baixo custo \aiflag{(Chmurzynski, 2024)}. Essa
problemática implica que, ainda que no longo prazo possa haver ganhos
com uma política de exportações, há no curto prazo uma disputa política
a respeito de qual consumo será inibido e qual classe, ou fração de
classe, arcará com tal custo.

Além disso, a busca de uma política de exportações tende a priorizar a
depreciação do câmbio, o que também é um elemento de restrição do
consumo. Nesse sentido, ela opera num sentido oposto de uma
sobrevalorização da taxa de câmbio, que permite manter padrões de
consumo elevados em troca de perda de competitividade nacional \aiflag{(Gala e
Libânio, 2011)}. Novamente, a problemática é menos técnica/econômica,
definir que um papel exportador é no longo prazo positivo a meta de
desenvolvimento econômico, e mais política, definir como a restrição do
consumo e seus supostos frutos futuros serão em cada caso repartidos
pelos grupos sociais.

Por fim, há de se ponderar o conjunto de argumentos voltados ao
desenvolvimento da infraestrutura urbana, além do estabelecimento de uma
indústria de saúde e defesa mais autossuficientes \citep{cndi_nova_2024}.
Em que pese a importância e a capacidade da melhoria da infraestrutura
urbana e da indústria de saúde em beneficiarem diretamente a qualidade
de vida da população, além do valor estratégico do setor de defesa e da
penetração do apelo ``nacionalista'' em parte da população, tais setores
apresentam limitações dentro do atual quadro fiscal brasileiro. Isso
pois investimentos em defesa e infraestrutura são altamente dependentes
da capacidade de gasto do Estado, de políticas públicas de longo prazo e
nem sempre resultam nos efeitos em cadeia desejados. Especialmente no
setor de defesa, expectativas de que a adoção de políticas de
desenvolvimento nacional resultem em ramificações e desenvolvimento de
atividades econômicas ``locais'' -- por exemplo via políticas de
Arranjos Produtivos Locais -- , tem tido resultados abaixo das
expectativas, mesmo que positivos, e dificuldades de continuidade em
cenários de mudança de governo \aiflag{(Ramalho e Conceição, 2024)}. São assim,
exemplo de como não basta o correto desenho de política pública, mas
também a necessidade de construção de um projeto de desenvolvimento
durável no tempo e resistente às alternâncias de poder democráticas.

\subsubsection{Oportunidades da reindustrialização: um novo pacto baseado na produtividade?}

O cerne dos argumentos acima é a relação entre a existência de um
projeto político definido e a característica do processo de
desenvolvimento econômico. Isso implica o estabelecimento de uma relação
entre classes, ou frações de classes, que seja funcional a manutenção de
tal projeto político. As formas históricas que esse pacto assume variam
no tempo em acordo com os interesses e o poder -- a capacidade de
implementar seus interesses como interesses comuns \aiflag{(Dobbin, 2005)} -- dos
diversos grupos sociais envolvidos.

No caso brasileiro o período de forte industrialização de 1967 até 1973
foi exemplo de um desses pactos. Como forma de sumarizar tal projeto,
poderia-se recorrer à frase do antigo ministro da economia Delfim Neto
de que é necessário ``fazer o bolo crescer, para depois dividi-lo''
(Folha de São Paulo, {[}s.d.{]}). Essa visão expressada é a de um
desenvolvimento econômico acelerado, ancorado em dívida externa e papel
ativo do Estado, mas que requer significativo aumento da
desigualdade\aiflag{(Barbosa, 2022; Hermann, 2015)}. É uma ótima demonstração da
característica de um pacto firmado em ambiente de autoritarismo e onde
predominava a conciliação dos interesses de elites econômicas nacionais,
internacionais e da burocracia estatal \citep{evans_dependent_1979}. A consequência da
desarticulação desse projeto, em meio a crise da dívida externa e
inflação, foi a articulação de um projeto de desenvolvimento
alternativo, baseado em abertura comercial, privatização e, acima de
tudo, estabilidade inflacionária\aiflag{(Giambiagi, 2015; Paula, de, 2022)}.

Entretanto, se bem sucedido no combate inflacionário, tal projeto não
evitou a continuidade da desindustrialização nacional. Se os governos
Lula {[}2003-2010{]} e Dilma {[}2011-2015{]} em parte conseguiram uma
parcial repactuação desse modelo, com a inclusão da expansão da
capacidade de consumo das classes populares entre os interesses
privilegiados, especialmente via valorização salarial \aiflag{(Cosenza, 2022;
Fonseca, Arend e Guerrero, 2020)}, tão pouco conseguiu reverter uma
tendência à desindustrialização relativa -- crescimento menor da
indústria que outros setores -- que se prolongou nas últimas duas
décadas.

Em outras palavras, em meio ao uso de recursos políticos e econômicos em
projetos de redistribuição de renda bem sucedidos, o papel do Estado de
ampliar a capacidade de acumulação ficou em segundo plano. O resultado
foi um crescimento econômico robusto, baseado no incremento da demanda
pela classes média emergente, no bônus demográfico e no período de
\emph{boom} das \emph{commodities}, mas de pouca melhoria na
produtividade, pouco reposicionamento nas cadeias globais de valor e
recorrentes pressões inflacionárias \aiflag{(Cosenza, 2022; Fonseca, Arend e
Guerrero, 2020)}.

Também argumentou-se que, historicamente, os projetos de desenvolvimento
baseados em significativa industrialização dependem de crescentes
volumes de capital, ação estatal e se articularam em torno de dois
eixos: a aceleração da capacidade de crescimento econômico baseada em
incrementos da produtividade e a oferta significativa de empregos. Esses
eixos dizem respeito diretamente a interesses de frações da elite
econômica, tipicamente uma fração industrialista, e das camadas
populares.

Após hiato entre 2016 e 2022 na existência de uma política industrial
articulada e aposta num modelo liberal, pautado na redução dos custos do
trabalho, enxugamento do estado e exploração das ``vantagens
competitivas'', o país voltou a enfatizar a necessidade de um
desenvolvimento econômico que incorpore significativa reindustrialização
como Nova Indústria Brasil. Entretanto, a reestruturação da atividade
industrial em curso desde meados do século XX, que resultou em processos
de desindustrialização em diversos países, reduziu de forma
significativa a capacidade de geração de empregos dos investimentos
industriais. Em resposta a isso e também a pressões particulares do caso
brasileiro, a nova política industrial iniciada em 2024 parte da
articulação do argumento pró-indústria com outros interesses e demandas
da sociedade, a saber: uma melhoria na qualidade do emprego oferecido, a
possibilidade de um modelo de desenvolvimento mais ambientalmente
sustentável, a melhoria da infraestrutura urbana utilizada pela
população e a autonomia nacional em setores estratégicos, como saúde e
defesa.

Entretanto, mesmo sendo significativo avanço, tal projeto de
reindustrialização ainda encontra obstáculos. Alguns dos maiores têm
relação com a significativa constrição fiscal -- em parte decorrente da
baixa produtividade -- e da pouca capacidade de agregação de apoio
popular diante da baixa capacidade de gerar empregos e de um ambiente
institucional hostil a conversão de processos de \emph{upgrade}
econômico em \emph{upgrade} social. Nesse cenário, a continuidade de
esforços de incremento da capacidade de consumo da população,
especialmente via valorização real do salário mínimo, parece significar
uma priorização do papel redistributivo do Estado, mas encontra
resistência nas pressões inflacionárias, de causas em parte internas e
em parte externas. A resposta institucionalizada a essa pressão
inflacionária é um aumento significativo da taxa de juros, que reprime a
possibilidade de desenvolvimento industrial. Acende sinal de alerta a
estagnação do setor industrial nos dados do PIB do primeiro trimestre de
2025, com especial preocupação para a retração de 1\% da indústria de
transformação \aiflag{(Agência de Notícias - IBGE, 2025)}.

Daí a importância de se atentar cada vez mais ao conjunto de demandas
populares que pode constituir base de apoio político a novos projetos de
desenvolvimento. Esse elemento participativo é fundamental para
constituir um modelo de desenvolvimento que concilie as demandas
específicas da atividade industrial com as demandas de uma sociedade
democrática. Dentre muitas possibilidades, destaca-se aqui o
significativo impacto e popularidade de projetos relativos à redução da
jornada de trabalho, em particular a forte rejeição popular da chamada
``escala 6x1" \aiflag{(Martins, 2024)}. Se os incrementos de produtividade do
trabalho, inerentes e necessários ao desenvolvimento industrial, são
contraproducentes a projetos baseados na oferta quantitativa de
empregos, eles são aparentemente mais compatíveis com projetos baseados
no retorno dessa produtividade não apenas na forma de salários, e
consumo, mas na forma da redução do tempo de trabalho. Para isso,
entretanto, é preciso constituir um ambiente institucional que permita a
captação dessa produtividade pelos trabalhadores e não apenas pelo
capital, o que vai na contramão de um conjunto de reformas realizadas e
propostas nos últimos anos para (des)regulação do trabalho e capital.

A articulação dessa e de outras demandas sociais enquanto objetivos
plausíveis, para a população, de um projeto de reindustrialização é
precondição da mesma. Elas atuariam potencialmente como fonte de capital
político para um desenvolvimento econômico pautado em atividades mais
sofisticadas e intensivas em capital e menos ancorado naquelas
intensivas em trabalho. Parece, assim, constituir um dos principais
desafios e oportunidades na busca da reindustrialização brasileira.

